% Options for packages loaded elsewhere
\PassOptionsToPackage{unicode}{hyperref}
\PassOptionsToPackage{hyphens}{url}
%
\documentclass[
  10pt,
  letterpaper,
]{article}
\usepackage{amsmath,amssymb}
\usepackage{iftex}
\ifPDFTeX
  \usepackage[T1]{fontenc}
  \usepackage[utf8]{inputenc}
  \usepackage{textcomp} % provide euro and other symbols
\else % if luatex or xetex
  \usepackage{unicode-math} % this also loads fontspec
  \defaultfontfeatures{Scale=MatchLowercase}
  \defaultfontfeatures[\rmfamily]{Ligatures=TeX,Scale=1}
\fi
\usepackage{lmodern}
\ifPDFTeX\else
  % xetex/luatex font selection
\fi
% Use upquote if available, for straight quotes in verbatim environments
\IfFileExists{upquote.sty}{\usepackage{upquote}}{}
\IfFileExists{microtype.sty}{% use microtype if available
  \usepackage[]{microtype}
  \UseMicrotypeSet[protrusion]{basicmath} % disable protrusion for tt fonts
}{}
\makeatletter
\@ifundefined{KOMAClassName}{% if non-KOMA class
  \IfFileExists{parskip.sty}{%
    \usepackage{parskip}
  }{% else
    \setlength{\parindent}{0pt}
    \setlength{\parskip}{6pt plus 2pt minus 1pt}}
}{% if KOMA class
  \KOMAoptions{parskip=half}}
\makeatother
\usepackage{xcolor}
\usepackage[margin=0.88in]{geometry}
\usepackage{longtable,booktabs,array}
\usepackage{calc} % for calculating minipage widths
% Correct order of tables after \paragraph or \subparagraph
\usepackage{etoolbox}
\makeatletter
\patchcmd\longtable{\par}{\if@noskipsec\mbox{}\fi\par}{}{}
\makeatother
% Allow footnotes in longtable head/foot
\IfFileExists{footnotehyper.sty}{\usepackage{footnotehyper}}{\usepackage{footnote}}
\makesavenoteenv{longtable}
\setlength{\emergencystretch}{3em} % prevent overfull lines
\providecommand{\tightlist}{%
  \setlength{\itemsep}{0pt}\setlength{\parskip}{0pt}}
\setcounter{secnumdepth}{5}
\usepackage{amsthm,mathtools}
\usepackage{enumitem}
\newtheorem{theorem}{Theorem}[section]
\newtheorem{lemma}[theorem]{Lemma}
\newtheorem{proposition}[theorem]{Proposition}
\newtheorem{corollary}[theorem]{Corollary}
\theoremstyle{definition}
\newtheorem{definition}[theorem]{Definition}
\theoremstyle{remark}
\newtheorem{remark}[theorem]{Remark}

\usepackage{mathrsfs}
\ifLuaTeX
  \usepackage{selnolig}  % disable illegal ligatures
\fi
\usepackage{bookmark}
\IfFileExists{xurl.sty}{\usepackage{xurl}}{} % add URL line breaks if available
\urlstyle{same}
\hypersetup{
  pdftitle={Transverse Plücker Arithmetic Holonomy for Catalan's Constant},
  hidelinks,
  pdfcreator={LaTeX via pandoc}}

\title{Transverse Plücker Arithmetic Holonomy for Catalan's Constant}
\usepackage{etoolbox}
\makeatletter
\providecommand{\subtitle}[1]{% add subtitle to \maketitle
  \apptocmd{\@title}{\par {\large #1 \par}}{}{}
}
\makeatother
\subtitle{An Audited Candidate Proof}
\author{}
\date{26 August 2026}

\begin{document}
\maketitle

{
\setcounter{tocdepth}{2}
\tableofcontents
}
\section{Abstract}\label{abstract}

Let \[
G=\beta(2)=\sum_{n\ge0}\frac{(-1)^n}{(2n+1)^2}
\] be Catalan's constant. This manuscript assembles a candidate proof of
\(G\notin\mathbf Q\) from a collection of independently audited
arithmetic, reflection, Hermite, denominator, and Archimedean modules.
The proof constructs a single rational rank-\(2D\) determinant line. At
\(p=2\), two transverse approximations produce a resource \[
\Omega_D=4D^2+4D-4\sum_{m=1}^D s_2(m),
\] while exact reflection/Hermite comparison and a conservative third
reflection reserve leave an effective coefficient \(\log2\). At odd
primes the descended rank-two source has denominator profile \((0,2)\),
giving cost \(3\). At infinity the actual conditional reflection channel
admits the quadratic fold \[
\phi_f(z)=\frac{64z(1-z)^2}{(1+z)^4},
\] with mixed energies \(T(r_0)=4\log2\), \(T(1)=8\log2\), and a
worst-case half-fold gain \[
I_*=\log2+\frac{6(\log2)^2}{\operatorname{arcosh}3}.
\] The adelic determinant inequality forced by rationality then reduces
to \[
2\log2-3+\frac{6(\log2)^2}{\operatorname{arcosh}3}\le0,
\] whereas a certified rational estimate proves that the left-hand side
is \(>27/1250\). The paper is deliberately labelled a candidate proof
pending independent expert verification of the assembled manuscript; the
purpose of the present version is to place every load-bearing transition
in one self-contained text.

\section{Main theorem and
architecture}\label{main-theorem-and-architecture}

\begin{theorem}[Main theorem -- candidate proof]\label{thm:main}
Catalan's constant
\[
G=\beta(2)=\sum_{n\ge0}\frac{(-1)^n}{(2n+1)^2}
\]
is irrational.
\end{theorem}

The proof is by contradiction. Assume throughout the global assembly
that \(G=a/b\in\mathbf Q\) in lowest terms. The argument is deliberately
modular. Sections A--E construct and analyze the finite arithmetic and
reflection source; Sections K-Finite and K-Global prove that the moving
Padé/Christoffel--Darboux and reflected Hermite realizations are
presentations of the same rational transverse determinant line. Section
G supplies the odd-prime denominator term. Sections H--J supply the
Archimedean fold and its optimal worst-case energy. The final two
sections reduce the product formula to a scalar inequality and
contradict it by certified rational bounds.

For traceability to the multi-party audit, the letters A--L are retained
in the section titles. They are not assumptions: the mathematical
content of the audited modules is reproduced below.

The manuscript uses \(v_2(2)=1\) and writes \(s_2(n)\) for the binary
digit sum. All asymptotic statements are as \(D\to\infty\).

\section{\texorpdfstring{Section A --- Exact \(2\)-adic transverse
arithmetic}{Section A --- Exact 2-adic transverse arithmetic}}\label{section-a-exact-2-adic-transverse-arithmetic}

\subsubsection{0. Conventions and
definitions}\label{conventions-and-definitions}

Normalize \(v_2(2)=1\), and write \(s_2(n)\) for the binary digit sum of
\(n\).

For \(k\ge 0\), put

\[
\left[\frac{k}{x}\right]
:=
\frac{k!}{x(x+1)\cdots(x+k)}.
\]

Let

\[
\Theta(x)
:=
-\sum_{k=0}^{\infty}
\left[\frac{k}{x}\right]
\left[\frac{k}{1-x}\right].
\]

This is the sign convention compatible with the Beukers continued
fraction and with

\[
\boxed{\Theta(x+1)+\Theta(x)=\frac{2}{x^2}.}
\]

Let \(q_n(x),p_n(x)\) be the Beukers denominator and numerator
polynomials, normalized by

\[
q_0(x)=1,\qquad q_1(x)=x^2-x+1,
\]

\[
p_0(x)=0,\qquad p_1(x)=1,
\]

and, for \(n\ge1\), by

\[
(n+1)^2u_{n+1}
=
\bigl(2n(n+1)+1-x+x^2\bigr)u_n
-
n^2u_{n-1}.
\]

Define

\[
E_n(x):=q_n(x)\Theta(x)-p_n(x).
\]

We shall also use the rising factorial

\[
(y)_r=y(y+1)\cdots(y+r-1),\qquad (y)_0=1.
\]

\begin{center}\rule{0.5\linewidth}{0.5pt}\end{center}

\subsection{A1. Beukers half-integer Padé
valuation}\label{a1.-beukers-half-integer-paduxe9-valuation}

For every odd integer \(a\) and every \(n\ge0\),

\[
\boxed{
v_2(E_n(a/2))=4n+2-2s_2(n),
}
\]

and

\[
\boxed{
v_2(q_n(a/2))=-4n+2s_2(n).
}
\]

\subsubsection{A1.1 Corrected remainder
expansion}\label{a1.1-corrected-remainder-expansion}

Beukers's Padé identity, in the above sign convention, is

\[
p_n(x)-q_n(x)\Theta(x)
=
(-1)^n
\sum_{k=n}^{\infty}
\binom{k}{n}
\left[\frac{k}{x}\right]
\left[\frac{k}{1-x}\right].
\]

Since \(E_n=q_n\Theta-p_n\),

\[
\boxed{
E_n(x)
=
(-1)^{n+1}
\sum_{k=n}^{\infty}
\binom{k}{n}
\left[\frac{k}{x}\right]
\left[\frac{k}{1-x}\right].
}
\tag{A1}
\]

This is the only sign correction needed in the original derivation.

\subsubsection{A1.2 Half-integer factorial
valuation}\label{a1.2-half-integer-factorial-valuation}

Let \(a\) be odd. Then

\[
\left[\frac{k}{a/2}\right]
=
\frac{k!}{\prod_{j=0}^{k}(a/2+j)}.
\]

Each factor \(a/2+j=(a+2j)/2\) has valuation \(-1\), because \(a+2j\) is
odd. Hence

\[
v_2\!\left(\prod_{j=0}^{k}(a/2+j)\right)=-(k+1).
\]

Legendre's formula gives

\[
v_2(k!)=k-s_2(k).
\]

Therefore

\[
\boxed{
v_2\!\left(\left[\frac{k}{a/2}\right]\right)
=
k-s_2(k)+(k+1)
=
2k+1-s_2(k).
}
\tag{A2}
\]

Since \(2-a\) is also odd, the same formula applies to
\(\left[k/(1-a/2)\right]\).

\subsubsection{A1.3 Unique minimum in the
remainder}\label{a1.3-unique-minimum-in-the-remainder}

The \(k\)-th summand in (A1) has valuation

\[
v_2\binom{k}{n}+4k+2-2s_2(k).
\]

At \(k=n\) this is

\[
4n+2-2s_2(n).
\]

Write \(k=n+r\), \(r>0\). The difference from the \(k=n\) valuation is

\[
4r
-
2\bigl(s_2(n+r)-s_2(n)\bigr)
+
v_2\binom{n+r}{n}.
\]

Binary digit-sum subadditivity gives

\[
s_2(n+r)-s_2(n)\le s_2(r)\le r.
\]

Hence the difference is at least

\[
4r-2r=2r>0.
\]

Thus the \(k=n\) term is uniquely of least \(2\)-adic valuation, so
cancellation is impossible:

\[
\boxed{
v_2(E_n(a/2))=4n+2-2s_2(n).
}
\tag{A3}
\]

\subsubsection{A1.4 Denominator
valuation}\label{a1.4-denominator-valuation}

The explicit Beukers formula is

\[
q_n(x)
=
\sum_{k=0}^{n}
\frac{(1-x)_{n-k}(x)_k^2}
{(n-k)!(k!)^2}.
\tag{A4}
\]

At \(x=a/2\), with \(a\) odd, the \(k\)-th summand has valuation

\[
V_k
=
-2n-2k+s_2(n-k)+2s_2(k).
\]

For \(k=n\),

\[
V_n=-4n+2s_2(n).
\]

Put \(j=n-k>0\). Then

\[
V_{n-j}-V_n
=
2j+s_2(j)+2s_2(n-j)-2s_2(n).
\]

Using

\[
s_2(n)\le s_2(j)+s_2(n-j),
\]

we obtain

\[
V_{n-j}-V_n
\ge
2j-s_2(j)>0.
\]

Thus the terminal term \(k=n\) is uniquely minimal, and

\[
\boxed{
v_2(q_n(a/2))=-4n+2s_2(n).
}
\tag{A5}
\]

\begin{center}\rule{0.5\linewidth}{0.5pt}\end{center}

\subsection{A2. Moving Rivoal--Zudilin
fibre}\label{a2.-moving-rivoalzudilin-fibre}

For \(m\ge1\), set

\[
x_m=\frac12-m,
\qquad
Q_m:=q_m(x_m),
\]

and define the alternating partial sum

\[
\sigma_m^{\mathrm{alt}}
:=
\sum_{j=0}^{m-1}\frac{(-1)^j}{(2j+1)^2}.
\]

Put

\[
\xi:=\Theta(1/2).
\]

Use the standard Rivoal--Zudilin normalization

\[
\boxed{
P_m
=
Q_m\sigma_m^{\mathrm{alt}}
+
\frac{(-1)^m}{8}\,p_m(x_m).
}
\tag{A6}
\]

This normalization begins with

\[
(Q_1,P_1)=\left(\frac74,\frac{13}{8}\right)
\]

and agrees with the usual Rivoal--Zudilin Catalan row.

\subsubsection{A2.1 Iteration of the functional
equation}\label{a2.1-iteration-of-the-functional-equation}

Let \(x_j=\frac12-j\). Since \(x_j+1=x_{j-1}\),

\[
\Theta(x_j)
=
\frac{2}{x_j^2}-\Theta(x_{j-1})
=
\frac{8}{(2j-1)^2}-\Theta(x_{j-1}).
\]

Iteration from \(x_0=1/2\) gives

\[
\boxed{
\Theta(x_m)
=
(-1)^m
\left(
\xi-8\sigma_m^{\mathrm{alt}}
\right).
}
\tag{A7}
\]

\subsubsection{A2.2 Exact moving error}\label{a2.2-exact-moving-error}

Using (A6) and (A7),

\[
\begin{aligned}
\frac{\xi}{8}-\frac{P_m}{Q_m}
&=
\frac{\xi}{8}
-
\sigma_m^{\mathrm{alt}}
-
\frac{(-1)^m}{8}\frac{p_m(x_m)}{q_m(x_m)}
\\
&=
\frac{(-1)^m}{8}
\left(
\Theta(x_m)-\frac{p_m(x_m)}{q_m(x_m)}
\right)
\\
&=
\frac{(-1)^m}{8Q_m}E_m(x_m).
\end{aligned}
\tag{A8}
\]

Now \(x_m=(1-2m)/2\) has odd numerator, so A1 applies:

\[
v_2(E_m(x_m))=4m+2-2s_2(m),
\]

\[
v_2(Q_m)=-4m+2s_2(m).
\]

Therefore

\[
\boxed{
v_2\!\left(
\frac{\xi}{8}-\frac{P_m}{Q_m}
\right)
=
8m-1-4s_2(m).
}
\tag{A9}
\]

Also

\[
\boxed{
v_2(Q_m)=-4m+2s_2(m).
}
\tag{A10}
\]

For later use, note that the \(k=0\) term of the series for \(\xi\) has
valuation \(2\), while all \(k\ge1\) terms have strictly larger
valuation. Hence

\[
v_2(\xi)=2,
\qquad
v_2(\xi/8)=-1.
\tag{A11}
\]

Since the error in (A9) has valuation \(>-1\) for \(m\ge1\),

\[
\boxed{
v_2(P_m)=-4m+2s_2(m)-1.
}
\tag{A12}
\]

\begin{center}\rule{0.5\linewidth}{0.5pt}\end{center}

\subsection{A3. Fixed modular fibre}\label{a3.-fixed-modular-fibre}

Define

\[
Y_0(t):=\sum_{k\ge0}q_k(1/2)t^k,
\qquad
Y_1(t):=\sum_{k\ge0}p_k(1/2)t^k,
\]

and

\[
A(z):=\sum_{n\ge0}A_nz^n,
\qquad
B(z):=\sum_{n\ge0}B_nz^n.
\]

Assume the stated quadratic pullback identities

\[
\boxed{
A(z)
=
(1-8z)\,
Y_0\!\left(16z(1-4z)\right),
}
\tag{A13}
\]

\[
\boxed{
16B(z)
=
(1-8z)\,
Y_1\!\left(16z(1-4z)\right).
}
\tag{A14}
\]

We prove

\[
\boxed{
v_2(A_n)=2s_2(n)\qquad(n\ge0),
}
\tag{A15}
\]

and

\[
\boxed{
v_2(B_n)=2s_2(n)-2\qquad(n\ge1).
}
\tag{A16}
\]

\subsubsection{A3.1 The fixed Beukers numerator
valuation}\label{a3.1-the-fixed-beukers-numerator-valuation}

At \(x=1/2\),

\[
p_n(1/2)=\xi q_n(1/2)-E_n(1/2).
\]

For \(n\ge1\),

\[
v_2(\xi q_n(1/2))
=
2-4n+2s_2(n),
\]

whereas

\[
v_2(E_n(1/2))
=
4n+2-2s_2(n).
\]

Their difference is

\[
8n-4s_2(n)>0.
\]

Thus \(\xi q_n(1/2)\) is uniquely of smaller valuation, and

\[
\boxed{
v_2(p_n(1/2))
=
-4n+2s_2(n)+2
\qquad(n\ge1).
}
\tag{A17}
\]

\subsubsection{\texorpdfstring{A3.2 Exact valuation of
\(A_n\)}{A3.2 Exact valuation of A\_n}}\label{a3.2-exact-valuation-of-a_n}

Put

\[
T(z)=16z(1-4z).
\]

In the unshifted part \(Y_0(T(z))\), the coefficient of \(z^n\) receives
contributions with

\[
k=n-j,
\qquad
0\le j\le \lfloor n/2\rfloor.
\]

The corresponding term is

\[
q_{n-j}(1/2)\,
16^{\,n-j}
(-4)^j
\binom{n-j}{j},
\]

of valuation

\[
2s_2(n-j)+2j+v_2\binom{n-j}{j}.
\tag{A18}
\]

For \(j=0\) this equals \(2s_2(n)\).

If \(j\ge2\), digit-sum subadditivity gives

\[
\begin{aligned}
&2s_2(n-j)+2j+v_2\binom{n-j}{j}
-2s_2(n)
\\
&\qquad\ge
2(j-s_2(j))
+
v_2\binom{n-j}{j}
>0.
\end{aligned}
\]

For \(j=1\), using

\[
s_2(n)=s_2(n-1)+1-v_2(n),
\]

the difference is

\[
v_2(n-1)+2v_2(n)>0
\]

whenever the \(j=1\) term exists.

Thus the \(j=0\) contribution is uniquely minimal in the unshifted
coefficient.

Every term coming from the factor \(-8z\) has valuation at least

\[
3+2s_2(n-1).
\]

Since

\[
s_2(n-1)=s_2(n)-1+v_2(n),
\]

this exceeds \(2s_2(n)\) by

\[
1+2v_2(n)>0.
\]

Hence the global minimum is unique and

\[
\boxed{v_2(A_n)=2s_2(n).}
\]

\subsubsection{\texorpdfstring{A3.3 Exact valuation of
\(B_n\)}{A3.3 Exact valuation of B\_n}}\label{a3.3-exact-valuation-of-b_n}

For \(k\ge1\), (A17) differs from the \(q_k(1/2)\) valuation by exactly
\(2\). The same coefficient-extraction argument therefore gives, for
\(n\ge1\),

\[
v_2(16B_n)=2s_2(n)+2.
\]

Consequently

\[
\boxed{
v_2(B_n)=2s_2(n)-2.
}
\]

\subsubsection{A3.4 Exact fixed-fibre
error}\label{a3.4-exact-fixed-fibre-error}

Because the present section uses

\[
E_k=q_k\Theta-p_k,
\]

subtracting \(\xi A(z)\) from \(16B(z)\) gives

\[
\boxed{
16B(z)-\xi A(z)
=
-(1-8z)
\sum_{k\ge0}
E_k(1/2)\,
T(z)^k.
}
\tag{A19}
\]

The minus sign in (A19) is required by the chosen definition of \(E_k\);
it does not affect any valuation.

Consider the coefficient of \(z^{2m}\), \(m\ge1\).

\paragraph{Unshifted terms}\label{unshifted-terms}

Write

\[
k=m+r,\qquad 0\le r\le m.
\]

Then the power of \((-4z)\) chosen from \((1-4z)^k\) is \(m-r\). The
term has valuation

\[
10m+6r+2
-
2s_2(m+r)
+
v_2\binom{m+r}{m-r}.
\]

At \(r=0\), this is

\[
10m+2-2s_2(m).
\tag{A20}
\]

For \(r>0\), the excess over (A20) is

\[
6r
+
2\bigl(s_2(m)-s_2(m+r)\bigr)
+
v_2\binom{m+r}{m-r}.
\]

Since

\[
s_2(m+r)-s_2(m)\le s_2(r)\le r,
\]

the excess is at least \(4r>0\).

\paragraph{\texorpdfstring{Terms from
\(-8z\)}{Terms from -8z}}\label{terms-from--8z}

Here \(0\le r\le m-1\), and the excess over (A20) is

\[
6r+1
+
2\bigl(s_2(m)-s_2(m+r)\bigr)
+
v_2\binom{m+r}{m-1-r},
\]

which is at least

\[
4r+1>0.
\]

Thus the \(r=0\) unshifted term is uniquely minimal. Hence

\[
\boxed{
v_2(16B_{2m}-\xi A_{2m})
=
10m+2-2s_2(m).
}
\tag{A21}
\]

Since

\[
v_2(A_{2m})
=
2s_2(2m)
=
2s_2(m),
\]

we obtain

\[
\boxed{
v_2\!\left(
\frac{\xi}{8}
-
\frac{2B_{2m}}{A_{2m}}
\right)
=
10m-1-4s_2(m).
}
\tag{A22}
\]

\begin{center}\rule{0.5\linewidth}{0.5pt}\end{center}

\subsection{A4. Transverse Plücker
determinant}\label{a4.-transverse-pluxfccker-determinant}

Define

\[
\rho_m:=\frac{4B_{2m}}{A_{2m}},
\qquad
\sigma_m:=\frac{2P_m}{Q_m},
\qquad m\ge1.
\]

It is useful to make the normalization step explicit. Before multiplying
by \(2\), the two approximations to the common \(2\)-adic quantity
\(\xi/8\) are

\[
\frac{P_m}{Q_m}
\quad\text{and}\quad
\frac{2B_{2m}}{A_{2m}}.
\]

Their error valuations are respectively

\[
8m-1-4s_2(m)
\]

and

\[
10m-1-4s_2(m).
\]

They are unequal for every \(m\ge1\). Therefore the strong ultrametric
equality gives

\[
v_2\!\left(
\frac{P_m}{Q_m}
-
\frac{2B_{2m}}{A_{2m}}
\right)
=
8m-1-4s_2(m).
\]

Multiplying the difference by \(2\),

\[
\boxed{
v_2(\sigma_m-\rho_m)
=
8m-4s_2(m).
}
\tag{A23}
\]

Now

\[
\sigma_m-\rho_m
=
\frac{
2(P_mA_{2m}-2B_{2m}Q_m)
}{
Q_mA_{2m}
}.
\]

Using

\[
v_2(Q_m)=-4m+2s_2(m)
\]

and

\[
v_2(A_{2m})=2s_2(m),
\]

equation (A23) is equivalent to

\[
\boxed{
v_2(P_mA_{2m}-2B_{2m}Q_m)
=
4m-1.
}
\tag{A24}
\]

Thus the transverse determinant statement is proved, not assumed.

\begin{center}\rule{0.5\linewidth}{0.5pt}\end{center}

\subsection{A5. Quadratic resource}\label{a5.-quadratic-resource}

Define

\[
\Omega_D
:=
\sum_{m=1}^{D}
v_2(\sigma_m-\rho_m).
\]

By (A23),

\[
\begin{aligned}
\Omega_D
&=
\sum_{m=1}^{D}
\bigl(8m-4s_2(m)\bigr)
\\
&=
8\frac{D(D+1)}2
-
4\sum_{m=1}^{D}s_2(m)
\\
&=
\boxed{
4D^2+4D-4S_D,
}
\end{aligned}
\]

where

\[
S_D:=\sum_{m=1}^{D}s_2(m).
\]

Since \(s_2(m)\le 1+\lfloor\log_2m\rfloor\),

\[
S_D=O(D\log D),
\]

and therefore

\[
\boxed{
\Omega_D=4D^2+O(D\log D).
}
\tag{A25}
\]

\begin{center}\rule{0.5\linewidth}{0.5pt}\end{center}

\subsection{A6. Independent small-row
check}\label{a6.-independent-small-row-check}

As a consistency check, one can compute the moving row directly from the
Beukers recurrence at \(x_m=1/2-m\), compute the fixed row directly from
the quadratic pullback, and use exact rational arithmetic.

For \(m=1,\ldots,8\), the resulting valuations are:

\begin{longtable}[]{@{}
  >{\raggedleft\arraybackslash}p{(\columnwidth - 8\tabcolsep) * \real{0.2000}}
  >{\raggedleft\arraybackslash}p{(\columnwidth - 8\tabcolsep) * \real{0.2000}}
  >{\raggedleft\arraybackslash}p{(\columnwidth - 8\tabcolsep) * \real{0.2000}}
  >{\raggedleft\arraybackslash}p{(\columnwidth - 8\tabcolsep) * \real{0.2000}}
  >{\raggedleft\arraybackslash}p{(\columnwidth - 8\tabcolsep) * \real{0.2000}}@{}}
\toprule\noalign{}
\begin{minipage}[b]{\linewidth}\raggedleft
\(m\)
\end{minipage} & \begin{minipage}[b]{\linewidth}\raggedleft
\(v_2(P_mA_{2m}-2B_{2m}Q_m)\)
\end{minipage} & \begin{minipage}[b]{\linewidth}\raggedleft
predicted \(4m-1\)
\end{minipage} & \begin{minipage}[b]{\linewidth}\raggedleft
\(v_2(\sigma_m-\rho_m)\)
\end{minipage} & \begin{minipage}[b]{\linewidth}\raggedleft
predicted \(8m-4s_2(m)\)
\end{minipage} \\
\midrule\noalign{}
\endhead
\bottomrule\noalign{}
\endlastfoot
1 & 3 & 3 & 4 & 4 \\
2 & 7 & 7 & 12 & 12 \\
3 & 11 & 11 & 16 & 16 \\
4 & 15 & 15 & 28 & 28 \\
5 & 19 & 19 & 32 & 32 \\
6 & 23 & 23 & 40 & 40 \\
7 & 27 & 27 & 44 & 44 \\
8 & 31 & 31 & 60 & 60 \\
\end{longtable}

This check is not used in the proof, but it catches the relevant
\(2/4/8\) normalization errors very effectively.

\begin{center}\rule{0.5\linewidth}{0.5pt}\end{center}

\section{Section D --- Reflection coinvariants and the rational
source}\label{section-d-reflection-coinvariants-and-the-rational-source}

\subsubsection{D1. Reflection module}\label{d1.-reflection-module}

Let \[
R_D=\mathbf Z[Y]_{<D}
\] viewed as a free \(\mathbf Z\)-module of rank \(D\). Set \[
L_D=R_DA\oplus R_DB\oplus R_DC.
\] Thus \(L_D\) is free of rank \(3D\), and \(T\) acts \(R_D\)-linearly.

\subsubsection{D2. Primitive anti-invariant summand and coinvariant
quotient}\label{d2.-primitive-anti-invariant-summand-and-coinvariant-quotient}

Let \[
N_D=R_D(B+C).
\] Since \[
(A,B,B+C)
\] is an integral basis of each coefficient slot (the change-of-basis
determinant from \((A,B,C)\) is \(1\)), one has the direct decomposition
\[
L_D=R_DA\oplus R_DB\oplus N_D.
\] Consequently \(N_D\) is primitive and saturated over \(\mathbf Z\),
including after localization at \(p=2\), and \[
W_D:=L_D/N_D
\] is torsion-free of rank \(2D\). Explicitly, \[
W_D\cong R_D[A]\oplus R_D[B],
\qquad [C]=-[B].
\]

Moreover \[
(T-1)A=0,\qquad (T-1)B=-(B+C),\qquad (T-1)C=-(B+C),
\] so \[
N_D=(T-1)L_D.
\] Hence \[
\boxed{W_D=L_D/(T-1)L_D}
\] is canonically the \textbf{coinvariant}, equivalently maximal
\(T\)-trivial, quotient. The induced action of \(T\) on \(W_D\) is the
identity.

This must not be confused integrally with the invariant submodule. In
fact \[
L_D^T=R_DA\oplus R_D(B-C),
\] and the natural map \[
L_D^T\longrightarrow W_D
\] sends \[
A\longmapsto[A],\qquad B-C\longmapsto2[B].
\] Therefore its cokernel is \[
(\mathbf Z/2\mathbf Z)^D.
\] After tensoring with \(\mathbf Q\), invariants and coinvariants are
canonically identified by averaging: \[
\boxed{
\jmath:W_D\otimes\mathbf Q\xrightarrow{\sim}L_D^T\otimes\mathbf Q,
\qquad
[v]\longmapsto\frac{v+Tv}{2}.
}
\] Thus \[
\jmath([A])=A,
\qquad
\jmath([B])=\frac{B-C}{2}.
\] The factor \(2^{-D}\) in this rational identification is the entire
invariant/coinvariant defect at \(p=2\); it is only \(O(D)\)
logarithmically.

\subsubsection{D3. Rationality-dependent invariant
basis}\label{d3.-rationality-dependent-invariant-basis}

Assume \[
G=\frac ab\in\mathbf Q,
\qquad \gcd(a,b)=1,
\qquad b>0.
\] Define \[
H_+=b(B-C)-aA,
\qquad
H_-=B+C.
\] Then \[
TH_+=H_+,
\qquad
TH_-=-H_-.
\] In \(W_D\), \[
[H_+]=2b[B]-a[A].
\] Hence, over \(\mathbf Q\), the change from the coinvariant basis
\(([A],[B])\) to \(([A],[H_+])\) has one-slot matrix \[
\begin{pmatrix}
1&-a\\
0&2b
\end{pmatrix}
\] and determinant \(2b\). Across \(D\) polynomial slots the determinant
is \[
(2b)^D.
\] Equivalently, the integral lattice generated by \([A],[H_+]\) has
index \((2b)^D\) in the lattice generated by \([A],[B]\). Since \(G\) is
fixed while \(D\to\infty\), this contributes only \[
D\log(2b)=O(D)
\] to any logarithmic determinant ledger.

\subsubsection{D4. Formal descent to the quotient
coordinate}\label{d4.-formal-descent-to-the-quotient-coordinate}

Let \[
s(X)=\sqrt{1-X}\in\mathbf Q[[X]],
\qquad s(0)=1.
\] Since \[
1-w(X)=\frac1{1-X},
\] one has, with the formal branch normalized at \(0\), \[
s(w(X))=s(X)^{-1}.
\] If \(Tf=f\), then \(f(w(X))=(1-X)f(X)=s(X)^2f(X)\), and therefore \[
(s f)(w(X))=s(w(X))f(w(X))=s(X)f(X).
\] Thus multiplication by \(s\) converts every \(T\)-invariant function
into an ordinary \(w\)-invariant function.

To identify the invariant formal-series ring, put \[
u=\frac{X}{2-X}.
\] Then \[
u(w(X))=-u(X),
\qquad
Y=-\frac{4u^2}{1-u^2}.
\] Since \(u\) is a formal coordinate at \(X=0\), and \(Y\) is a formal
coordinate in \(u^2\), \[
\mathbf Q[[X]]^w
=\mathbf Q[[u]]^{u\mapsto-u}
=\mathbf Q[[u^2]]
=\mathbf Q[[Y]].
\] Therefore the two invariant combinations descend unconditionally: \[
\boxed{
P(Y)=s(X)A(X)\in\mathbf Q[[Y]],
}
\] \[
\boxed{
R_0(Y)=s(X)(B-C)(X)\in\mathbf Q[[Y]].
}
\]

Under the rationality hypothesis define \[
K(Y)=s(X)(B-C-GA)(X).
\] Then \[
\boxed{K=R_0-GP\in\mathbf Q[[Y]].}
\] The basis change \[
(P,R_0)\longmapsto(P,K=R_0-GP)
\] is the same global \(\mathbf Q\)-linear change at every place, with
matrix \[
\begin{pmatrix}
1&-G\\
0&1
\end{pmatrix}
\] and determinant \(1\). It need not preserve the standard
\(\mathbf Z_p\)-lattice at primes dividing \(b\), but it has exactly
zero determinant-line cost at every place.

The averaging/descent map from the coinvariant quotient satisfies \[
[A]\longmapsto P,
\qquad
[B]\longmapsto \frac{R_0}{2},
\] and, under \(G=a/b\), \[
[H_+]\longmapsto bK.
\] Thus the integral factor \(2\) separating invariants from
coinvariants is explicit and cannot be hidden in the descent.

\subsubsection{D5. The global rational
source}\label{d5.-the-global-rational-source}

To avoid assuming a zero estimate at this stage, define the source first
as the abstract rational vector space \[
\boxed{
\mathcal E_D^{\mathrm{abs}}
=
\mathbf Q[Y]_{<D}e_P
\oplus
\mathbf Q[Y]_{<D}e_{R_0},
}
\] which has dimension exactly \(2D\), together with the global
evaluation map \[
\operatorname{ev}_D:\mathcal E_D^{\mathrm{abs}}\longrightarrow\mathbf Q[[Y]],
\qquad
f e_P+g e_{R_0}\longmapsto fP+gR_0.
\]

If one writes \[
\mathbf Q[Y]_{<D}P\oplus\mathbf Q[Y]_{<D}R_0
\] as a direct sum \textbf{inside} \(\mathbf Q[[Y]]\), injectivity of
\(\operatorname{ev}_D\) is additionally required. That is not a
consequence of the reflection quotient alone; it is proved separately in
Section E. Accordingly, injectivity is deferred to that section.

Under \(G\in\mathbf Q\), the abstract basis may globally be changed to
\[
\mathbf Q[Y]_{<D}e_P
\oplus
\mathbf Q[Y]_{<D}e_K,
\qquad e_K=e_{R_0}-G e_P,
\] with determinant \(1\). At the Archimedean place these are the \(D\)
base directions and \(D\) fold directions used later.

\section{Shared local expansion
lemma}\label{shared-local-expansion-lemma}

Place this immediately after the construction of the descended source
\((P,R_0)\), before the Hermite, denominator, and fold sections.

Put \[
t=\frac Y4,\qquad z=\frac Y{64},
\] and \[
\mathscr Q=4t(1-t)D_t^2+4(1-2t)D_t-1.
\] The reflected companion satisfies \[
\mathscr Q\,R_0(4t)=\frac1{2\sqrt{1-t}},
\qquad R_0(0)=0,
\] and \[
P(Y)={}_2F_1\!\left(\frac12,\frac12;1;\frac Y4\right).
\]

Write \[
R_0(4t)=\sum_{n\ge0}c_nt^n.
\] After division by \(4\), coefficient extraction gives \[
(n+1)^2c_{n+1}-\left(n+\frac12\right)^2c_n
=
\frac18\frac{\binom{2n}{n}}{4^n}.
\tag{S0.1}
\]

Let \[
a_n=\left(\frac{(1/2)_n}{n!}\right)^2
=\frac{\binom{2n}{n}^2}{16^n},
\] so \[
(n+1)^2a_{n+1}=\left(n+\frac12\right)^2a_n.
\] Define \[
S_0=0,\qquad
S_n=\sum_{m=0}^{n-1}
\frac{4^m}{(2m+1)^2\binom{2m}{m}}.
\] Then \[
S_{n+1}-S_n
=
\frac{4^n}{(2n+1)^2\binom{2n}{n}}.
\]

The sequence \[
c_n=\frac12a_nS_n
\] satisfies (S0.1), since \[
\begin{aligned}
&(n+1)^2c_{n+1}-\left(n+\frac12\right)^2c_n\\
&=
\frac12\left(n+\frac12\right)^2a_n(S_{n+1}-S_n)
=
\frac18\frac{\binom{2n}{n}}{4^n}.
\end{aligned}
\] It also has \(c_0=0\), so uniqueness of the first-order recurrence
gives \[
R_0(4t)=
\frac12\sum_{n\ge1}
\frac{\binom{2n}{n}^2}{16^n}S_nt^n.
\]

Putting \(t=16z\), \[
\boxed{
R_0(64z)=
\frac12\sum_{n\ge1}\binom{2n}{n}^2
\left(
\sum_{m=0}^{n-1}
\frac{4^m}{(2m+1)^2\binom{2m}{m}}
\right)z^n.
}
\tag{S0.2}
\] Also \[
\boxed{
P(64z)=\sum_{n\ge0}\binom{2n}{n}^2z^n.
}
\tag{S0.3}
\]

Under \(G\in\mathbf Q\), \(K=R_0-GP\), hence \[
\boxed{
K(64z)=
\sum_{n\ge0}\binom{2n}{n}^2
\left[
-G+
\frac12\sum_{m=0}^{n-1}
\frac{4^m}{(2m+1)^2\binom{2m}{m}}
\right]z^n.
}
\tag{S0.4}
\]

This is exactly the local expansion used in Sections E and H. Section H
proves the inverse-central-binomial identity converting (S0.4) to the
tail form.

For the final manuscript, the inhomogeneous equation for \(R_0\) should
be proved at the same point as the descent/reflection differential
equations; audited Section K already uses it in the exact fixed-fibre
bridge.

\section{Section E --- Reflected Hermite and Hankel
normality}\label{section-e-reflected-hermite-and-hankel-normality}

\subsubsection{E1. Exact coefficients and their 2-adic
size}\label{e1.-exact-coefficients-and-their-2-adic-size}

Set \[
z=\frac Y{64}.
\] From \[
\left(\frac{(1/2)_n}{n!}\right)^2
=\frac1{16^n}\binom{2n}{n}^2,
\] the hypergeometric expansion of \(P\) gives \[
P(64z)=\sum_{n\ge0}p_nz^n,
\qquad
p_n=\binom{2n}{n}^2.
\]

Using the displayed local expansion of \(K\) and \(R_0=K+GP\), the
\(-G\)-term cancels coefficientwise and \[
R_0(64z)=\sum_{n\ge1}r_nz^n,
\] where \[
r_n=\frac12\binom{2n}{n}^2S_n,
\qquad
S_n=\sum_{m=0}^{n-1}
\frac{4^m}{(2m+1)^2\binom{2m}{m}}.
\] (The constant term is zero.)

For \(m\ge1\), Kummer's identity \[
v_2\binom{2m}{m}=s_2(m)
\] gives \[
v_2\!\left(
\frac{4^m}{(2m+1)^2\binom{2m}{m}}
\right)
=2m-s_2(m)\ge1.
\] The \(m=0\) summand is \(1\). Hence \[
S_n\in\mathbf Z_2^\times,
\qquad
S_n\equiv1\pmod2.
\] Consequently \[
v_2(r_n)=2s_2(n)-1\ge1,
\] so \(r_n\in2\mathbf Z_2\), and \[
\frac{r_n}{2}\equiv1\pmod2
\iff s_2(n)=1.
\]

\subsubsection{E2. Ratio parity theorem}\label{e2.-ratio-parity-theorem}

Define \[
\beta(z)=\frac{R_0(64z)}{P(64z)}
=\sum_{n\ge1}b_nz^n,
\qquad b_0=0.
\] Since \(p_0=1\), coefficient comparison in \[
\left(\sum_{k\ge0}p_kz^k\right)
\left(\sum_{j\ge0}b_jz^j\right)
=\sum_{n\ge0}r_nz^n
\] gives, for \(n\ge1\), \[
b_n=r_n-\sum_{k=1}^{n-1}p_kb_{n-k}.
\tag{E2.1}
\] For every \(k\ge1\), \[
v_2(p_k)=2s_2(k)\ge2,
\qquad p_k\in4\mathbf Z_2.
\] We prove by induction on \(n\) that \(b_n\in2\mathbf Z_2\) and \[
b_n\equiv r_n\pmod8.
\tag{E2.2}
\] For \(n=1\), (E2.1) has no correction term, so the claim is
immediate. If it holds below \(n\), then every summand \(p_kb_{n-k}\) in
(E2.1) lies in \[
4\mathbf Z_2\cdot2\mathbf Z_2=8\mathbf Z_2.
\] Since \(r_n\in2\mathbf Z_2\), both assertions follow.

Dividing (E2.2) by \(2\) gives the stronger congruence \[
\frac{b_n}{2}\equiv\frac{r_n}{2}\pmod4,
\] and therefore, in particular, \[
\frac{b_n}{2}\equiv\frac{r_n}{2}\pmod2.
\] By E1, \[
\boxed{
\frac{b_n}{2}\equiv1\pmod2
\iff s_2(n)=1
\iff n\text{ is a positive power of }2.
}
\]

\subsubsection{E3. Binary Hankel
determinant}\label{e3.-binary-hankel-determinant}

Let \[
M_D^{(z)}=(b_{D+i-j})_{0\le i,j<D}.
\] Reverse the \(D\) columns. The reversal permutation has sign \[
(-1)^{D(D-1)/2},
\] and the resulting matrix is \[
H_D=(b_{i+j+1})_{0\le i,j<D}.
\] Thus \[
\det M_D^{(z)}=(-1)^{D(D-1)/2}\det H_D,
\] so column reversal does not affect the \(2\)-adic valuation.

By E2 every \(b_n\in2\mathbf Z_2\). Hence \[
H_D=2C_D,
\qquad C_D\in M_D(\mathbf Z_2),
\] and modulo \(2\), \[
C_D\equiv\mathcal A_D=(a_{i+j+1})_{0\le i,j<D},
\qquad
a_n=\mathbf1_{\{n=2^r\}}.
\] Also define \[
\mathcal B_D=(a_{i+j+2})_{0\le i,j<D}.
\] All determinant identities in the rest of this section are in
\(\mathbf F_2\).

For \(m\ge1\), \[
a_{2m}=a_m,
\qquad
a_{2m+1}=0,
\] while \(a_1=1\).

\paragraph{\texorpdfstring{Even size for
\(\mathcal A\)}{Even size for \textbackslash mathcal A}}\label{even-size-for-mathcal-a}

For \(D=2h\), order rows and columns by even indices first and odd
indices second. Then \[
\mathcal A_{2h}\sim
\begin{pmatrix}
E_h&\mathcal A_h\\
\mathcal A_h&0
\end{pmatrix},
\] where \(E_h\) has a single \(1\) in its \((0,0)\)-entry. Swapping the
two column blocks gives a block-upper-triangular matrix with diagonal
blocks \(\mathcal A_h,\mathcal A_h\). Therefore \[
\det\mathcal A_{2h}=(\det\mathcal A_h)^2.
\tag{E3.1}
\]

\paragraph{\texorpdfstring{Odd size for
\(\mathcal A\)}{Odd size for \textbackslash mathcal A}}\label{odd-size-for-mathcal-a}

For \(D=2h+1\) with \(h\ge1\), parity ordering followed by separating
the even index \(0\) gives \[
\mathcal A_{2h+1}\sim
\begin{pmatrix}
1&0&u^T\\
0&0&\mathcal B_h\\
u&\mathcal B_h&0
\end{pmatrix}
\] for some \(u\in\mathbf F_2^h\). Adding suitable multiples of the
first column to the last \(h\) columns, and then suitable multiples of
the first row to the last \(h\) rows, clears \(u\) without changing the
determinant. Hence \[
\det\mathcal A_{2h+1}=(\det\mathcal B_h)^2.
\tag{E3.2}
\]

\paragraph{\texorpdfstring{The \(\mathcal B\)
recurrences}{The \textbackslash mathcal B recurrences}}\label{the-mathcal-b-recurrences}

The same parity ordering gives directly \[
\mathcal B_{2h}\sim
\begin{pmatrix}
\mathcal A_h&0\\
0&\mathcal B_h
\end{pmatrix},
\] so \[
\det\mathcal B_{2h}
=\det\mathcal A_h\,\det\mathcal B_h,
\tag{E3.3}
\] and \[
\mathcal B_{2h+1}\sim
\begin{pmatrix}
\mathcal A_{h+1}&0\\
0&\mathcal B_h
\end{pmatrix},
\] so \[
\det\mathcal B_{2h+1}
=\det\mathcal A_{h+1}\,\det\mathcal B_h.
\tag{E3.4}
\]

Now \[
\mathcal A_1=(a_1)=(1),
\qquad
\mathcal B_1=(a_2)=(1).
\] A simultaneous strong induction on \(D\), using (E3.1)--(E3.4),
yields \[
\det\mathcal A_D=
\det\mathcal B_D=1
\qquad(D\ge1)
\] in \(\mathbf F_2\).

Therefore \(\det C_D\in\mathbf Z_2^\times\), and \[
\det H_D=2^D\det C_D.
\] Consequently \[
\boxed{v_2(\det M_D^{(z)})=D.}
\]

\subsubsection{E4. Hermite independence}\label{e4.-hermite-independence}

Let \(J_{2D}\) denote the first \(2D\) coefficient map in the
\(z\)-coordinate. Because \[
P(64z)=1+O(z),
\] multiplication by \(P(64z)^{-1}\) induces on \(J_{2D}\) a
lower-triangular target automorphism with diagonal entries \(1\). Thus
it preserves rank and determinant nonvanishing.

After dividing the source by \(P\), and using the \(\mathbf Q\)-basis \[
1,z,\ldots,z^{D-1},
\qquad
\beta,z\beta,\ldots,z^{D-1}\beta,
\] the coefficient matrix for orders \(0,1,\ldots,2D-1\) has block form
\[
\begin{pmatrix}
I_D&T_D\\
0&M_D^{(z)}
\end{pmatrix},
\] where the lower-right block is exactly \[
(M_D^{(z)})_{i,j}=b_{D+i-j}.
\] Hence the full \(2D\times2D\) coefficient determinant equals
\(\det M_D^{(z)}\), and is nonzero by E3.

The original polynomial source is written in powers of \(Y\), whereas
the displayed matrix uses powers of \(z=Y/64\). The two polynomial bases
differ by a nonzero diagonal rational rescaling, so this change does not
affect linear independence. (Its \(2\)-adic determinant must of course
be charged separately whenever one tracks the determinant in the
unscaled \(Y\)-coordinate.)

Therefore the first \(2D\) coefficient evaluations on \(\mathcal E_D\)
are linearly independent for every \(D\).

Under the rationality hypothesis, Section D gives the unipotent basis
change \[
(P,R_0)\longmapsto(P,K=R_0-GP),
\] so the \(2D\) source directions split as exactly \(D\) directions of
type \(P\) and \(D\) directions of type \(K\).

\section{Section C --- Exact reflection Smith
mass}\label{section-c-exact-reflection-smith-mass}

\subsubsection{C1. Reflection matrix and exact entry
valuations}\label{c1.-reflection-matrix-and-exact-entry-valuations}

For \(0\le m\le r<D\), \[
U_{r,m}
=
(-1)^m\frac{(m+\frac12)_{2(r-m)}}{(2(r-m))!},
\] and \(U_{r,m}=0\) for \(m>r\).

Put \(n=r-m\). Then \[
(m+\tfrac12)_{2n}
=2^{-2n}\prod_{j=0}^{2n-1}(2m+1+2j),
\] and every factor in the product is odd. Hence \[
v_2\bigl((m+\tfrac12)_{2n}\bigr)=-2n.
\] Legendre's formula gives \[
v_2((2n)!)=2n-s_2(2n)=2n-s_2(n).
\] Therefore \[
\boxed{
v_2(U_{r,m})=-4(r-m)+s_2(r-m).
}
\] This also holds for \(r=m\), since both sides are \(0\).

The matrix is lower triangular and \[
U_{m,m}=(-1)^m.
\] Consequently \[
\boxed{
\det U_D=(-1)^{D(D-1)/2}\in\{\pm1\}.
}
\]

\subsubsection{C2. Universal lower bound for determinant terms and
minors}\label{c2.-universal-lower-bound-for-determinant-terms-and-minors}

Consider a nonzero determinant term in any \(k\times k\) minor. Write
its matched row/column pairs as \((r_i,m_{\sigma(i)})\), and put \[
d_i=r_i-m_{\sigma(i)}\ge0.
\] Let \[
I=\{i:d_i>0\},\qquad q=|I|,\qquad S=\sum_i d_i=\sum_{i\in I}d_i.
\] The \(q\) positive source indices \(m_{\sigma(i)}\) are distinct
elements of \(\{0,\ldots,D-1\}\), and so are the \(q\) positive target
indices \(r_i\). Therefore \[
\sum_{i\in I}r_i
\le (D-q)+(D-q+1)+\cdots+(D-1),
\] while \[
\sum_{i\in I}m_{\sigma(i)}
\ge 0+1+\cdots+(q-1).
\] Subtracting gives \[
\boxed{S\le q(D-q).}
\]

By C1, the valuation of the determinant term is \[
-4S+\sum_i s_2(d_i).
\] Every positive integer has binary digit sum at least \(1\), hence \[
\sum_i s_2(d_i)\ge q.
\] Thus every nonzero determinant term satisfies \[
v_2(\text{term})
\ge -4q(D-q)+q.
\] Set \[
h=\lfloor D/2\rfloor,
\qquad c=\lceil D/2\rceil,
\qquad
R_D:=4hc-h.
\] Equivalently, \[
\boxed{R_D=D^2-\lceil D/2\rceil.}
\] For every integer \(q\), \[
4q(D-q)-q\le R_D.
\] Indeed, if \(D=2h\), the difference is \[
R_D-[4q(D-q)-q]
=4(q-h)^2+(q-h)\ge0,
\] and if \(D=2h+1\), it is \[
R_D-[4q(D-q)-q]
=4(q-h)^2-3(q-h)\ge0.
\] (The displayed quadratic is nonnegative for every integer \(q-h\).)
Hence every determinant term has valuation at least \(-R_D\), and by the
ultrametric inequality every minor satisfies \[
\boxed{
v_2(\det M)\ge -R_D.
}
\] This is a sharp \emph{universal} lower bound over minors of all
sizes. It is not, in general, the exact minimum for each fixed size
\(k\).

\subsubsection{C3. A minor attaining the
bound}\label{c3.-a-minor-attaining-the-bound}

The case \(D=1\) is immediate: \(U_1=[1]\) and \(R_1=0\). Assume
\(D\ge2\), so \(h\ge1\).

Take the \(h\times h\) minor \[
M=(U_{c+i,j})_{0\le i,j<h},
\] using rows \(c,c+1,\ldots,D-1\) and columns \(0,1,\ldots,h-1\). Put
\[
d_{ij}=c+i-j.
\] All \(d_{ij}>0\). Define \[
\widetilde M_{ij}=2^{4d_{ij}-1}M_{ij}.
\] By C1, \[
v_2(\widetilde M_{ij})=s_2(d_{ij})-1\ge0,
\] so \(\widetilde M\in M_h(\mathbf Z_2)\). Modulo \(2\), \[
\overline{\widetilde M}_{ij}
=
\begin{cases}
1,&d_{ij}\text{ is a power of }2,\\
0,&\text{otherwise}.
\end{cases}
\] (The sign \((-1)^j\) and all odd units reduce to \(1\) in
\(\mathbf F_2\).)

Moreover the normalization is a genuine row/column scaling: \[
\widetilde M
=
\operatorname{diag}(2^{4(c+i)-1})_{i=0}^{h-1}
\,M\,
\operatorname{diag}(2^{-4j})_{j=0}^{h-1}.
\] Hence \[
v_2(\det\widetilde M)
=
v_2(\det M)+h(4c-1).
\] It remains to show that \(\det\widetilde M\) is odd.

Let \(a_n\in\mathbf F_2\) be the indicator that \(n\ge1\) is a power of
\(2\). Reverse the order of the columns. If \(D=2h\), the reduced matrix
becomes \[
A_h=(a_{i+j+1})_{0\le i,j<h};
\] if \(D=2h+1\), it becomes \[
B_h=(a_{i+j+2})_{0\le i,j<h}.
\] We prove simultaneously that every \(A_n\) and every \(B_n\) is
nonsingular over \(\mathbf F_2\).

The identities \[
a_{2m}=a_m\quad(m\ge1),
\qquad
a_{2m+1}=0\quad(m\ge1),
\qquad a_1=1
\] are immediate.

For \(B_n\), reorder rows and columns by parity of their indices. With
\[
e=\lceil n/2\rceil,
\qquad o=\lfloor n/2\rfloor,
\] one obtains the block diagonal matrix \[
B_n\sim A_e\oplus B_o.
\] Thus \(B_n\) is nonsingular once the smaller matrices on the right
are.

For \(A_n\), the same parity reordering gives \[
A_n\sim
\begin{pmatrix}
E&C\\
C^T&0
\end{pmatrix},
\] where \(E\) has a single \(1\) in its \((0,0)\)-entry and zeros
elsewhere, and \[
C_{p,q}=a_{p+q+1}.
\] If \(n=2m\), then \(C=A_m\). If \(A_m\) is nonsingular, a vector in
the kernel has even block zero from the lower equation and then odd
block zero from the upper equation. Hence \(A_{2m}\) is nonsingular.

If \(n=2m+1\), then \(C\) is \((m+1)\times m\). Its first \(m\) rows
form \(A_m\), so \(C\) has full column rank. Its last \(m\) rows form
\(B_m\). Suppose \((x,y)\) lies in the kernel. Then \[
C^T x=0,
\qquad
x_0e_0+Cy=0.
\] If \(x_0=0\), the first equation restricted to the last \(m\)
coordinates gives \(B_m^Tx_{1:}=0\), hence \(x=0\), and then \(y=0\). If
\(x_0=1\), the second equation says \(Cy=e_0\); deleting its first row
gives \(B_my=0\), hence \(y=0\), contradiction. Thus the kernel is
trivial.

Starting from \(A_1=B_1=[1]\), strong induction proves that all
\(A_n,B_n\) are nonsingular. Therefore \[
\det\widetilde M\in\mathbf Z_2^\times,
\qquad
v_2(\det\widetilde M)=0.
\] Consequently \[
v_2(\det M)
=-h(4c-1)
=-4hc+h
=-R_D.
\] So the universal lower bound from C2 is attained.

\subsubsection{C4. Smith exponents and exact
mass}\label{c4.-smith-exponents-and-exact-mass}

For \(0\le k\le D\), let \[
\delta_k
=
\min\{v_2(\det N):N\text{ is a nonzero }k\times k\text{ minor of }U_D\},
\] with \(\delta_0=0\). Smith theory over the DVR \(\mathbf Z_2\) gives
\[
\delta_k=e_1+\cdots+e_k.
\] By C2, \(\delta_k\ge -R_D\) for every \(k\). By C3, the \(h\times h\)
attaining minor gives \(\delta_h=-R_D\). Therefore \[
\min_k\delta_k=-R_D.
\] Since \(e_1\le\cdots\le e_D\), the minimum partial sum is exactly the
sum of the negative Smith exponents. Hence \[
\boxed{
-\sum_{e_i<0}e_i=R_D=D^2-\lceil D/2\rceil.
}
\]

From C1, \[
\det U_D=(-1)^{D(D-1)/2},
\] so \[
\sum_i e_i=v_2(\det U_D)=0.
\] Therefore \[
\boxed{
\sum_{e_i>0}e_i=R_D=D^2-\lceil D/2\rceil.
}
\]

\section{Section B --- Christoffel--Darboux
no-backflow}\label{section-b-christoffeldarboux-no-backflow}

The following local theorem is proved:

\[
\boxed{
\text{Intrinsic selected-mode Christoffel–Darboux transport is }
2\text{-adically unimodular}.
}
\]

For strictly increasing pivots \(1\le n_0<\cdots<n_{D-1}\), \[
T_{ij}
=
\mathbf1_{j\le i}
(\sigma_{n_j}-\rho_{n_j})
q_{n_j}(1/2-n_i)q_{n_j}(1/2)
\] satisfies \[
\boxed{T_D\in GL_D(\mathbf Z_2)}
\] and therefore \[
\boxed{v_2(\det T_D)=0.}
\]

The proof consists of the uniform valuation \[
v_2q_n(1/2-m)=-4n+2s_2(n),
\] the projective depth from Section A \[
v_2(\sigma_n-\rho_n)=8n-4s_2(n),
\] and the exact compensation \[
v_2\!\left(
(\sigma_n-\rho_n)
q_n(1/2-m)q_n(1/2)
\right)=0.
\]

The Christoffel--Darboux identity \[
(n+1)^2W_n(x,y)
=
(x-y)(x+y-1)
\sum_{k=0}^nq_k(x)q_k(y)
\] proves the required triangularity because the row indexed by \(n_i\)
contains only modes \(k\le n_i\).

If literal \(W_n\) rather than the normalized CD sum is used, row \(i\)
is additionally multiplied by \[
\frac{n_i^2}{(n_i+1)^2}.
\] For consecutive pivots \(N,\ldots,N+D-1\), these factors telescope to
\[
\left(\frac{N}{N+D}\right)^2,
\] so their \(2\)-adic determinant cost is \(O(\log D)\) when
\(N=O(D)\). More generally, \(O(D)\)-range pivots give
\(O(D\log D)=o(D^2)\).

Thus the CD transport itself cannot pay back a positive proportion of
the \(D^2\)-scale transverse \(2\)-adic resource.

A global compatibility statement is additionally required to ensure that
the same saturated integral source/quotient lattice survives subsequent
pullback, Hermite, and basis changes. This compatibility is proved
explicitly in Section K-Finite; the present CD theorem is the local
input to that construction.

\section{Section K-Finite --- Common determinant line and finite
compatibility}\label{section-k-finite-common-determinant-line-and-finite-compatibility}

\subsubsection{Audited and completed
version}\label{audited-and-completed-version}

\subsection{K1. General determinant-line
identity}\label{k1.-general-determinant-line-identity}

Let \[
\psi:E\longrightarrow F
\] be a rational linear isomorphism. At a place \(v\), let \(S_v\) and
\(T_v\) be source and target basis matrices relative to fixed rational
reference bases. Then \[
\boxed{
\log|\det(T_v^{-1}\psi S_v)|_v
=
\log|\det\psi|_v
+
\log|\det S_v|_v
-
\log|\det T_v|_v.
}
\tag{K1.1}
\]

Thus:

\begin{enumerate}
\def\labelenumi{\arabic{enumi}.}
\tightlist
\item
  different local bases and lattices may be used at different places;
\item
  every local source and target basis determinant must be charged at
  that place;
\item
  only a genuinely global rational basis change may be cancelled by the
  product formula;
\item
  unrelated local changes of coordinates do not cancel automatically.
\end{enumerate}

\begin{center}\rule{0.5\linewidth}{0.5pt}\end{center}

\subsection{K2. The single rational source and the exact fixed-fibre
bridge}\label{k2.-the-single-rational-source-and-the-exact-fixed-fibre-bridge}

Section D supplies \[
\boxed{
\mathcal E_D^{\rm abs}
=
\mathbf Q[Y]_{<D}e_P
\oplus
\mathbf Q[Y]_{<D}e_{R_0}.
}
\tag{K2.1}
\]

Its evaluation map is \[
f e_P+g e_{R_0}\longmapsto fP+gR_0.
\]

Section E proves that the first \(2D\) coefficient functionals are
independent.

Under \(G\in\mathbf Q\), put \[
e_K=e_{R_0}-Ge_P.
\] Then \[
(e_P,e_{R_0})\longleftrightarrow(e_P,e_K)
\] is a global rational unipotent basis change of determinant exactly
\(1\).

The remaining invariant/coinvariant and integral rationality-basis
discrepancies in D have logarithmic size only \(O(D)\), hence
\(o(D^2)\).

\subsubsection{K2.1 Exact identification of Section A's fixed modular
fibre}\label{k2.1-exact-identification-of-section-as-fixed-modular-fibre}

Write Section A's modular series as \[
A_{\rm mod}(\zeta),\qquad B_{\rm mod}(\zeta).
\]

Define \[
Y_0(t)=\sum_{n\ge0}q_n(1/2)t^n,
\qquad
Y_1(t)=\sum_{n\ge0}p_n(1/2)t^n.
\]

Section A gives \[
A_{\rm mod}(\zeta)
=(1-8\zeta)Y_0(16\zeta(1-4\zeta)),
\tag{K2.2}
\] \[
16B_{\rm mod}(\zeta)
=(1-8\zeta)Y_1(16\zeta(1-4\zeta)).
\tag{K2.3}
\]

Put \[
t=16\zeta(1-4\zeta),\qquad g=1-8\zeta.
\] Then \[
1-t=g^2.
\tag{K2.4}
\]

At \(x=1/2\), the Padé recurrence gives \[
\mathscr L_B
=
t(t-1)^2D_t^2
+(3t-1)(t-1)D_t
+t-\frac34
\] with \[
\mathscr L_BY_0=0,\qquad
\mathscr L_BY_1=1.
\tag{K2.5}
\]

For \[
V(t)=\sqrt{1-t}\,Y(t),
\] direct substitution gives \[
\mathscr L_B((1-t)^{-1/2}V)
=
\frac{\sqrt{1-t}}4\,\mathscr QV,
\tag{K2.6}
\] where \[
\mathscr Q
=
4t(1-t)D_t^2
+4(1-2t)D_t
-1.
\tag{K2.7}
\]

For \[
V_0=\sqrt{1-t}\,Y_0
\] we have \[
\mathscr QV_0=0,\qquad
V_0(0)=1,\qquad
V_0'(0)=\frac14.
\]

Since \[
P(Y)
={}_2F_1\!\left(\frac12,\frac12;1;\frac Y4\right),
\] the function \(P(4t)\) satisfies the same equation and initial data.
Therefore \[
\boxed{
\sqrt{1-t}\,Y_0(t)=P(4t).
}
\tag{K2.8}
\]

Similarly, \[
V_1=\sqrt{1-t}\,Y_1
\] satisfies \[
\mathscr QV_1=\frac4{\sqrt{1-t}},
\qquad
V_1(0)=0,\quad V_1'(0)=1.
\tag{K2.9}
\]

The factored equation for \(R_0\), after \(Y=4t\), gives \[
\mathscr Q R_0(4t)=\frac1{2\sqrt{1-t}}.
\tag{K2.10}
\]

Since Section E gives \[
R_0(Y)=\frac Y{32}+O(Y^2),
\] the function \(8R_0(4t)\) has the same initial data as \(V_1\). Hence
\[
\boxed{
\frac18\sqrt{1-t}\,Y_1(t)=R_0(4t).
}
\tag{K2.11}
\]

Consequently \[
\boxed{
A_{\rm mod}(\zeta)
=
P\!\left(64\zeta(1-4\zeta)\right),
}
\tag{K2.12}
\] \[
\boxed{
2B_{\rm mod}(\zeta)
=
R_0\!\left(64\zeta(1-4\zeta)\right).
}
\tag{K2.13}
\]

Thus Section A's fixed projective fibre is literally a quadratic
pullback of the same \((P,R_0)\)-source.

Set \[
z=\zeta(1-4\zeta).
\] Because \[
z=\zeta+O(\zeta^2),
\] composition by this map on every finite jet space is triangular with
diagonal \(1\). Its determinant is therefore \(1\).

The only non-unit coordinate scale is \[
Y=64z.
\]

\begin{center}\rule{0.5\linewidth}{0.5pt}\end{center}

\subsection{\texorpdfstring{K3. The common saturated transverse quotient
at
\(p=2\)}{K3. The common saturated transverse quotient at p=2}}\label{k3.-the-common-saturated-transverse-quotient-at-p2}

Work over \(\mathbf Q_2\) in \[
z=\frac Y{64}.
\]

Define \[
\beta(z)=\frac{R_0(64z)}{P(64z)}.
\]

Since \[
P(64z)=1+O(z),
\] division by \(P(64z)\) is a lower-triangular \(2\)-adic unit
operation on jets.

Using the source basis \[
1,z,\ldots,z^{D-1},
\qquad
\beta,z\beta,\ldots,z^{D-1}\beta,
\] Section E gives \[
\begin{pmatrix}
I_D&T_D\\
0&M_D^{(z)}
\end{pmatrix}.
\tag{K3.1}
\]

Hence \[
\Lambda_D^{\rm pol}
=
\langle1,z,\ldots,z^{D-1}\rangle_{\mathbf Z_2}
\] is primitive and saturated, and \[
\boxed{
\mathcal Q_D
=
\Lambda_D^{(z)}/\Lambda_D^{\rm pol}
}
\tag{K3.2}
\] is an explicit rank-\(D\) transverse lattice with basis represented
by \[
\beta,z\beta,\ldots,z^{D-1}\beta.
\]

Choose the consecutive CD pivots \[
1,2,\ldots,D.
\] The CD sum in row \(i\) contains the modes \(0,\ldots,i\). The
selected modes are \(1,\ldots,i\), so the only nonselected mode is \[
q_0=1,
\] which lies in the polynomial block.

Thus Section B's common-complement hypothesis is satisfied by the
explicit saturated polynomial complement. No separate polynomial-witness
conjecture is needed.

\begin{center}\rule{0.5\linewidth}{0.5pt}\end{center}

\subsection{K4. Exact CD
factorization}\label{k4.-exact-cd-factorization}

For \(i,j=1,\ldots,D\), define \[
h_j=4^j(j!)^2.
\tag{K4.1}
\]

Section B proves \[
h_jq_j\!\left(\frac12-i\right)\in\mathbf Z_2^\times,
\qquad
h_jq_j\!\left(\frac12\right)\in\mathbf Z_2^\times.
\tag{K4.2}
\]

Put \[
r_{ij}=h_jq_j\!\left(\frac12-i\right),
\qquad
c_j=h_jq_j\!\left(\frac12\right),
\] and \[
L_D=(\mathbf1_{j\le i}r_{ij}),
\qquad
C_D^{(0)}=\operatorname{diag}(c_1,\ldots,c_D),
\qquad
H_D=\operatorname{diag}(h_1,\ldots,h_D).
\tag{K4.3}
\]

Then \[
L_D,C_D^{(0)}\in GL_D(\mathbf Z_2).
\tag{K4.4}
\]

The unweighted selected-mode matrix is \[
\boxed{
\mathcal C_D
=
L_DC_D^{(0)}H_D^{-2}.
}
\tag{K4.5}
\]

Since \[
v_2(h_j)=4j-2s_2(j)
\] and \[
v_2(\Delta_j)=8j-4s_2(j),
\] we can write \[
\Delta_j=u_jh_j^2,
\qquad
u_j\in\mathbf Z_2^\times.
\tag{K4.6}
\]

Hence \[
\boxed{
\mathcal C_D
\operatorname{diag}(\Delta_1,\ldots,\Delta_D)
=
L_DC_D^{(0)}
\operatorname{diag}(u_1,\ldots,u_D)
\in GL_D(\mathbf Z_2).
}
\tag{K4.7}
\]

This is the exact no-backflow statement on the explicit quotient
\(\mathcal Q_D\).

Put \[
A_D=v_2(\det H_D).
\] Then \[
A_D
=
2D(D+1)-2S_D
\tag{K4.8}
\] and therefore \[
\boxed{
\Omega_D=2A_D.
}
\tag{K4.9}
\]

The literal \(W_n\) row normalization contributes only \(O(\log D)\) for
these consecutive pivots.

\begin{center}\rule{0.5\linewidth}{0.5pt}\end{center}

\subsection{K5. Exact reflected Hermite
factorization}\label{k5.-exact-reflected-hermite-factorization}

After column reversal, Section E's scaled transverse Hermite matrix is
\[
(b_{i+j+1})_{0\le i,j<D}.
\]

Write \[
(b_{i+j+1})=2V_D,
\qquad
V_D\in GL_D(\mathbf Z_2).
\tag{K5.1}
\]

Here \(V_D\) is unrelated to Section C's reflection matrix.

Set \[
A_{64}
=
\operatorname{diag}
(1,64^{-1},\ldots,64^{-(D-1)}).
\tag{K5.2}
\]

Returning to the unscaled coordinate \(Y=64z\), the transverse matrix is
exactly \[
\boxed{
\mathcal H_D
=
\frac1{32}A_{64}V_DA_{64}.
}
\tag{K5.3}
\]

Put \[
B_D=-v_2(\det A_{64})
=3D(D-1).
\tag{K5.4}
\]

Therefore \[
v_2(\det\mathcal H_D)
=
-5D-2B_D
=
-6D^2+D,
\] and \[
\boxed{
H_D^{\rm ref}=6D^2-D.
}
\tag{K5.5}
\]

The entire quadratic \(Y/64\) cost is visible in (K5.3), and must not be
charged again.

\begin{center}\rule{0.5\linewidth}{0.5pt}\end{center}

\subsection{K6. K-COMPAT: explicit common determinant-line
comparison}\label{k6.-k-compat-explicit-common-determinant-line-comparison}

For the Hermite presentation define \[
S_H=A_{64}^{-1}V_D^{-1},
\qquad
T_H=A_{64}.
\tag{K6.1}
\]

Then \[
\boxed{
T_H^{-1}\mathcal H_DS_H
=
\frac1{32}I_D.
}
\tag{K6.2}
\]

For the CD presentation define \[
S_C
=
\frac1{32}H_D(C_D^{(0)})^{-1}L_D^{-1},
\tag{K6.3}
\] \[
T_C
=
L_DH_D^{-1}L_D^{-1}.
\tag{K6.4}
\]

Then \[
\boxed{
T_C^{-1}\mathcal C_DS_C
=
\frac1{32}I_D.
}
\tag{K6.5}
\]

Set \[
S=S_HS_C^{-1},
\qquad
T=T_HT_C^{-1}.
\tag{K6.6}
\]

It follows that \[
\boxed{
\mathcal C_D=T^{-1}\mathcal H_DS.
}
\tag{K6.7}
\]

Thus the CD and reflected Hermite matrices are explicit local
presentations of the same rational transverse isomorphism.

\subsubsection{K6.1 Source transition}\label{k6.1-source-transition}

We have \[
v_2(\det S_H)=B_D,
\] and \[
v_2(\det S_C)=A_D-5D.
\]

Therefore \[
\begin{aligned}
v_2(\det S)
&=B_D-A_D+5D\\
&=\boxed{D^2+2S_D}.
\end{aligned}
\tag{K6.8}
\]

Hence \[
\boxed{
v_2(\det S)
=
R_D+O(D\log D).
}
\tag{K6.9}
\]

\subsubsection{K6.2 Target transition}\label{k6.2-target-transition}

Similarly, \[
v_2(\det T_C)=-A_D,
\qquad
v_2(\det T_H)=-B_D,
\] so \[
\boxed{
v_2(\det T)
=
-D^2+5D-2S_D.
}
\tag{K6.10}
\]

Therefore the adverse target mass is \[
\boxed{
-v_2(\det T)
=
R_D+O(D\log D).
}
\tag{K6.11}
\]

These are the two actual quadratic compatibility costs. They are not
asserted to be literally Section C's reflection matrix.

\subsubsection{K6.3 Reconciliation
identity}\label{k6.3-reconciliation-identity}

From K1 and (K6.7), \[
v_2(\det\mathcal C_D)-v_2(\det\mathcal H_D)
=
v_2(\det S)-v_2(\det T).
\]

Hence \[
\boxed{
H_D^{\rm ref}-\Omega_D
=
2D^2-5D+4S_D.
}
\tag{K6.12}
\]

Subtracting two reflection masses, \[
\boxed{
H_D^{\rm ref}-\Omega_D-2R_D
=
4S_D-5D+2\lceil D/2\rceil
=
O(D\log D).
}
\tag{K6.13}
\]

This is now a consequence of the explicit determinant-line comparison,
not merely a numerical coincidence.

Thus \textbf{K-COMPAT is proved}.

\begin{center}\rule{0.5\linewidth}{0.5pt}\end{center}

\subsection{\texorpdfstring{K7. The \(p=2\) reflection
budget}{K7. The p=2 reflection budget}}\label{k7.-the-p2-reflection-budget}

Section A supplies \[
\Omega_D\log2
=
4\log2\,D^2+o(D^2).
\]

K6 accounts for two actual compatibility transitions with total adverse
mass \[
2R_D\log2+o(D^2).
\]

Retain one additional complete mass \[
R_D\log2
\] as conservative reserve for any eliminated/free complement
normalization not used in the transverse comparison.

Therefore \[
\boxed{
\text{total reflection-scale charge}
\le
3R_D\log2+o(D^2).
}
\tag{K7.1}
\]

Since \[
R_D=D^2+O(D),
\] \[
\Omega_D-3R_D
=
D^2+o(D^2),
\] and hence \[
\boxed{
C_{2,\mathrm{eff}}=\log2.
}
\tag{K7.2}
\]

No fourth quadratic \(2\)-adic transition appears.

\begin{center}\rule{0.5\linewidth}{0.5pt}\end{center}

\section{Section G --- Odd-prime denominator
theorem}\label{section-g-odd-prime-denominator-theorem}

\subsubsection{G1. The depth-zero
direction}\label{g1.-the-depth-zero-direction}

Since \[
p_n=\binom{2n}{n}^2\in\mathbf Z,
\] the coefficients of \(P(64z)\) are integral.

Passing from \(z\) to \(Y=64z\) introduces only powers of \(2\). Hence
at every odd prime \(p\), \[
P(Y)
\] has denominator depth \(0\).

Thus the first source direction has odd-prime profile \[
\boxed{\sigma_1=0.}
\]

\subsubsection{\texorpdfstring{G2. Exact odd-prime denominator theorem
for
\(R_0\)}{G2. Exact odd-prime denominator theorem for R\_0}}\label{g2.-exact-odd-prime-denominator-theorem-for-r_0}

Put \[
d_n=\operatorname{lcm}(1,\dots,n).
\]

We prove directly from (G.1) that \[
\boxed{
d_n^2r_n\in\mathbf Z_{(2)}
}
\] for every \(n\), where \(\mathbf Z_{(2)}\) denotes the rationals
integral at every odd prime.

Fix an odd prime \(p\), and put \[
e=e_p(n):=v_p(d_n)
=\max\{j:p^j\le n\}.
\]

For \(0\le m<n\), write \[
a=v_p(2m+1),
\] and let \[
c_p(t)
=
v_p\binom{2t}{t}.
\] By Kummer's theorem, \(c_p(t)\) is the number of carries occurring
when \(t+t\) is added in base \(p\).

The \(m\)-th summand of \(r_n\) has \(p\)-adic valuation \[
2c_p(n)-2a-c_p(m),
\] because \(2\) and \(4\) are \(p\)-adic units. Thus after
multiplication by \(d_n^2\) its valuation is \[
2e+2c_p(n)-2a-c_p(m).
\tag{G.3}
\]

It remains to prove that (G.3) is nonnegative.

Since \[
p^a\mid 2m+1,
\] we have \[
m\equiv\frac{p^a-1}{2}\pmod{p^a}.
\] Therefore the lowest \(a\) base-\(p\) digits of \(m\) are all \[
\frac{p-1}{2}.
\] Adding \(m+m\) produces no carry in any of these lowest \(a\)
positions.

Since \[
m<n<p^{e+1},
\] there are at most \(e+1-a\) remaining positions in which a carry can
occur. Hence \[
c_p(m)\le e+1-a.
\tag{G.4}
\]

Also \[
2m+1<2p^{e+1}<p^{e+2},
\] because \(p\ge3\), and therefore \[
a\le e+1.
\]

If \[
a\le e-1,
\] then by (G.4) \[
2a+c_p(m)
\le
2a+e+1-a
=
e+a+1
\le2e.
\] Thus (G.3) is nonnegative.

Suppose next that \[
a=e.
\] Then \(c_p(m)\le1\). If \(c_p(m)=0\), the claim is immediate. If
\(c_p(m)=1\), the only possible carry is in the top base-\(p\) digit.
Writing \[
m
=
d\,p^e+\frac{p^e-1}{2},
\] that carry implies \[
2d\ge p.
\] Consequently \[
m>\frac{p^{e+1}}2.
\] Since \(n>m\), \[
n>\frac{p^{e+1}}2,
\] so \(n+n\) has a carry out of its top digit and therefore \[
c_p(n)\ge1.
\] Hence \[
2e+2c_p(n)\ge2e+2>2e+1
\ge2a+c_p(m).
\]

Finally suppose \[
a=e+1.
\] Because \(0\le m<p^{e+1}\), \[
m=\frac{p^{e+1}-1}{2}.
\] There are no carries in \(m+m\), so \[
c_p(m)=0.
\] Again \(n>m\) implies \[
n>\frac{p^{e+1}}2,
\] and hence \[
c_p(n)\ge1.
\] Therefore \[
2e+2c_p(n)
\ge
2e+2
=
2a.
\]

All cases give \[
2e+2c_p(n)-2a-c_p(m)\ge0.
\] Thus each summand of \(d_n^2r_n\) is \(p\)-integral, and so \[
\boxed{
d_n^2r_n\in\mathbf Z_p
\qquad(p\text{ odd}).
}
\tag{G.5}
\]

Since \(p\) was arbitrary, \[
\boxed{
d_n^2r_n\in\mathbf Z_{(2)}.
}
\]

\paragraph{Sharpness}\label{sharpness}

Take \(n=p\) for an odd prime \(p\). Then \[
v_p\binom{2p}{p}=0.
\]

For \[
m_0=\frac{p-1}{2},
\] one has \[
v_p(2m_0+1)=1,
\qquad
v_p\binom{2m_0}{m_0}=0.
\] Hence the \(m_0\)-summand in (G.1) has valuation exactly \[
-2.
\]

For every other \(0\le m<p\), \[
v_p(2m+1)=0,
\] while \[
v_p\binom{2m}{m}\le1.
\] Thus every other summand has valuation at least \(-1\).

The valuation \(-2\) is therefore attained uniquely and cannot cancel.
Hence \[
\boxed{
v_p(r_p)=-2
\qquad(p\text{ odd prime}).
}
\tag{G.6}
\]

Thus the second source direction has genuine odd-prime depth exactly
\(2\): \[
\boxed{\sigma_2=2.}
\]

Because \(Y=64z\) differs from \(z\) only by a power of \(2\),
(G.5)--(G.6) are unchanged in the unscaled \(Y\)-coordinate.

Consequently the reflected rank-two source has the exact odd-prime
denominator profile \[
\boxed{(0,2).}
\]

\subsubsection{G3. Compatibility with the reflection quotient and
rationality
basis}\label{g3.-compatibility-with-the-reflection-quotient-and-rationality-basis}

Section D gives \[
W_D
=
R_D[A]\oplus R_D[B]
\] as the integral coinvariant quotient.

The invariant lattice maps to \(W_D\) with cokernel \[
(\mathbf Z/2\mathbf Z)^D.
\] Therefore for every odd prime \(p\), invariants and coinvariants
define the same \(\mathbf Z_p\)-lattice.

Under \[
G=\frac ab,
\qquad
\gcd(a,b)=1,
\] Section D gives \[
[H_+]=2b[B]-a[A].
\] The corresponding one-slot transition matrix has determinant \(2b\),
so across \(D\) polynomial slots its determinant is \[
(2b)^D.
\] Its logarithmic contribution is \[
D\log(2b)=O(D)=o(D^2).
\]

Equivalently, after descent, \[
[H_+]\longmapsto bK,
\qquad
bK=bR_0-aP.
\] Since \(P\) has depth \(0\) and \(R_0\) has depth \(2\), \[
bK
\] has depth at most \(2\). At every odd prime not dividing the fixed
integer \(b\), (G.6) shows that the depth-\(2\) layer remains sharp. The
finitely many primes dividing \(b\) affect only the \(O(D)\) fixed-basis
term and do not alter the \(D^2\)-coefficient.

Hence the rationality-adapted rank-two source has asymptotic odd-prime
profile \[
\boxed{(0,2).}
\]

\subsubsection{G4. CDT denominator
rearrangement}\label{g4.-cdt-denominator-rearrangement}

Represent the denominator type by two identical columns \[
\mathbf b
=
\begin{pmatrix}
0&0\\
1&1
\end{pmatrix}.
\] The row sums are \[
(\sigma_1,\sigma_2)=(0,2).
\]

For \(m=2\), (G.2) gives \[
\tau
=
\frac{1\cdot0+3\cdot2}{4}
=
\boxed{\frac32}.
\]

In the polynomial-source filtration used here, there are \(mD=2D\)
source directions. The standard CDT/Bost \(D^2\)-normalization converts
the rearranged denominator rate by \[
\frac{m^2}{2}\tau.
\] For \(m=2\), \[
\frac{m^2}{2}\tau
=
2\tau
=
2\cdot\frac32
=
\boxed3.
\]

Equivalently, the relevant filtration has \(2D\) successive orders,
whose total size is \[
0+1+\cdots+(2D-1)
=
2D^2-D,
\] so the normalized leading scale carries the factor \(m^2/2=2\).

Thus the odd-prime denominator term charged in the global slopes ledger
is \[
\boxed{
3D^2+o(D^2).
}
\]

This is the CDT denominator \textbf{upper cost}. No assertion that every
actual determinant has precisely this denominator is required.

\subsubsection{\texorpdfstring{G5. Exclusion of
\(p=2\)}{G5. Exclusion of p=2}}\label{g5.-exclusion-of-p2}

Define the odd part of the lcm by \[
d_n^{\rm odd}
=
\frac{d_n}{2^{\lfloor\log_2n\rfloor}}.
\] Then \[
\log d_n^{\rm odd}
=
\log d_n
-
\lfloor\log_2n\rfloor\log2.
\]

Since \[
\log d_n=\psi(n)=n+o(n),
\] removing the prime \(2\) changes the logarithm at order \(n\) by only
\[
O(\log n).
\] Across \(O(D)\) coefficient orders this contributes \[
O(D\log D)=o(D^2).
\]

Therefore deleting \(p=2\) does not alter the leading odd-prime
denominator coefficient.

All \(2\)-adic reflection, Smith, averaging, Hermite, and no-backflow
terms are to be charged separately in Sections A--F. There is no double
counting.

\section{Section H --- Analytic continuation of the actual fold
channel}\label{section-h-analytic-continuation-of-the-actual-fold-channel}

\subsubsection{H0. Local coefficient
lemma}\label{h0.-local-coefficient-lemma}

Put \[
z=\frac{Y}{64}.
\] The conditional \(K\)-germ at \(Y=0\) has the exact expansion \[
\boxed{
K(64z)
=
\sum_{n\ge0}\binom{2n}{n}^2
\left[
-G+
\frac12\sum_{m=0}^{n-1}
\frac{4^m}{(2m+1)^2\binom{2m}{m}}
\right]z^n .
}
\tag{H0.1}
\]

This is the only extra local datum not contained in Section D.

We will also need the classical inverse-central-binomial identity \[
\boxed{
\sum_{m=0}^{\infty}
\frac{4^m}{(2m+1)^2\binom{2m}{m}}
=2G.
}
\tag{H0.2}
\]

For completeness, (H0.2) has a short direct proof. Since \[
\frac{1}{(2m+1)\binom{2m}{m}}
=
\int_0^1 t^m(1-t)^m\,dt
\] and \[
\frac1{2m+1}=\int_0^1x^{2m}\,dx,
\] monotone convergence gives \[
\begin{aligned}
\sum_{m\ge0}
\frac{4^m}{(2m+1)^2\binom{2m}{m}}
&=
\int_0^1\!\!\int_0^1
\frac{dt\,dx}{1-4t(1-t)x^2}\\
&=
\int_0^1
\frac{\arcsin x}{x\sqrt{1-x^2}}\,dx\\
&=
\int_0^{\pi/2}\frac{\theta}{\sin\theta}\,d\theta\\
&=
-2\int_0^{\pi/4}\log(\tan u)\,du\\
&=2G.
\end{aligned}
\]

Subtracting (H0.2) from the finite sum in (H0.1) therefore gives the
exact tail form \[
\boxed{
K(64z)
=
-\frac12\sum_{n\ge0}
\binom{2n}{n}^2
\left(
\sum_{m=n}^{\infty}
\frac{4^m}{(2m+1)^2\binom{2m}{m}}
\right)z^n .
}
\tag{H0.3}
\]

\begin{center}\rule{0.5\linewidth}{0.5pt}\end{center}

\subsubsection{H1. Differential
equation}\label{h1.-differential-equation}

Let \[
c_n=[z^n]K(64z).
\] Using \[
\frac{\binom{2n+2}{n+1}}{\binom{2n}{n}}
=
\frac{2(2n+1)}{n+1}
\] together with the first-difference relation for the partial sums in
(H0.1), direct elimination gives \[
\boxed{
(n+1)(n+2)^2c_{n+2}
-
4(n+1)(8n^2+18n+11)c_{n+1}
+
32(2n+1)^3c_n
=0.
}
\tag{H1.1}
\]

Equivalently, if \[
K(Y)=\sum_{n\ge0}k_nY^n,
\qquad c_n=64^nk_n,
\] then \[
128(n+1)(n+2)^2k_{n+2}
-
8(n+1)(8n^2+18n+11)k_{n+1}
+
(2n+1)^3k_n
=0.
\tag{H1.2}
\]

Coefficient extraction is equivalent to \[
\boxed{
8Y(Y-4)^2K'''
+4(Y-4)(9Y-16)K''
+2(13Y-44)K'
+K=0.
}
\tag{H1.3}
\]

As an independent algebraic check, the operator factors exactly as \[
\boxed{
\bigl[-2(Y-4)D_Y-1\bigr]
\circ
\bigl[
4Y(4-Y)D_Y^2+(16-8Y)D_Y-1
\bigr].
}
\tag{H1.4}
\]

Hence the only finite singular points of the scalar equation are \[
Y=0,\qquad Y=4.
\]

\begin{center}\rule{0.5\linewidth}{0.5pt}\end{center}

\subsubsection{\texorpdfstring{H2. Local exponents at
\(Y=4\)}{H2. Local exponents at Y=4}}\label{h2.-local-exponents-at-y4}

Put \(t=Y-4\) and insert \[
K\sim t^\rho.
\] The lowest-order terms are \[
32t^2K'''+80tK''+16K',
\] so the indicial coefficient is \[
\begin{aligned}
&32\rho(\rho-1)(\rho-2)
+80\rho(\rho-1)
+16\rho\\
&\qquad=
16\rho^2(2\rho-1).
\end{aligned}
\] After removing the nonzero scalar \(16\), the indicial polynomial is
\[
\boxed{\rho^2(2\rho-1)}.
\] Thus the local exponents are \[
\boxed{0,\quad0,\quad\frac12}.
\]

Because the exponent \(0\) is repeated, a logarithmic Frobenius solution
is possible a priori.

\begin{center}\rule{0.5\linewidth}{0.5pt}\end{center}

\subsubsection{H3. Exact coefficient asymptotic and exclusion of the
logarithm}\label{h3.-exact-coefficient-asymptotic-and-exclusion-of-the-logarithm}

Write \[
q_m=
\frac{4^m}{(2m+1)^2\binom{2m}{m}}.
\] Stirling's formula gives \[
\binom{2m}{m}
=
\frac{4^m}{\sqrt{\pi m}}
\left(1+O(m^{-1})\right),
\] and hence \[
q_m
=
\frac{\sqrt\pi}{4}\,m^{-3/2}
\left(1+O(m^{-1})\right).
\] Therefore \[
\sum_{m=n}^{\infty}q_m
=
\frac{\sqrt\pi}{2}\,n^{-1/2}
\left(1+O(n^{-1})\right).
\tag{H3.1}
\]

Also \[
\binom{2n}{n}^2
=
\frac{16^n}{\pi n}
\left(1+O(n^{-1})\right).
\tag{H3.2}
\] Combining (H0.3), (H3.1), and (H3.2) gives \[
\boxed{
[z^n]K(64z)
=
-\frac{16^n}{4\sqrt\pi\,n^{3/2}}
\left(1+O(n^{-1})\right).
}
\tag{H3.3}
\]

In particular, \[
\limsup_{n\to\infty}|[z^n]K(64z)|^{1/n}=16,
\] so the radius of convergence in \(z\) is \(1/16\), equivalently the
radius in the unscaled \(Y\)-coordinate is \(4\).

By H1 the only nonzero finite singular point of the differential
equation is \(Y=4\). Thus \(Y=4\) is the unique possible dominant finite
singularity of the germ centered at \(Y=0\).

Now suppose the chosen \(K\)-solution had a nonzero logarithmic
connection coefficient at \(Y=4\). A Frobenius logarithmic term has the
form \[
c\,H(Y)\log(1-Y/4),
\qquad H(4)\ne0,\quad c\ne0.
\] In the \(z\)-coordinate \(Y=64z\), singularity analysis gives \[
[z^n]\bigl(cH(64z)\log(1-16z)\bigr)
=
-cH(4)\frac{16^n}{n}
+O\!\left(\frac{16^n}{n^2}\right).
\tag{H3.4}
\] Because \(Y=4\) is the unique dominant finite singularity, there is
no second same-modulus singularity from which such an \(n^{-1}\)
contribution could cancel.

But (H3.3) is of exact order \[
16^n n^{-3/2},
\] which is strictly smaller than \(16^n n^{-1}\). Hence \(c=0\).

Therefore the logarithmic connection coefficient vanishes, and in a
neighborhood of \(Y=4\) \[
\boxed{
K(Y)
=
A_0(Y)
+
\sqrt{1-Y/4}\,A_1(Y),
}
\tag{H3.5}
\] with \(A_0,A_1\) holomorphic there.

This proves the required connection statement exactly.

\begin{center}\rule{0.5\linewidth}{0.5pt}\end{center}

\subsubsection{H4. Quadratic
desingularization}\label{h4.-quadratic-desingularization}

Set \[
\boxed{
F(U)=4U-U^2=4-(U-2)^2.
}
\] Then \[
1-\frac{F(U)}4
=
\frac{(U-2)^2}{4}.
\] Consequently the square-root term in (H3.5) becomes holomorphic after
composition with \(F\).

Define, initially as a germ at \(U=0\), \[
g(U)=K(F(U)).
\] Using the chain rule in (H1.3) gives exactly \[
\boxed{
U(U-2)(U-4)g'''
+
2(3U^2-12U+8)g''
+
7(U-2)g'
+
g
=0.
}
\tag{H4.1}
\]

The finite singular points of this pulled-back scalar equation are \[
U=0,\qquad U=2,\qquad U=4.
\]

For the particular solution \(g\):

\begin{itemize}
\tightlist
\item
  \(U=0\) is removable because \(F(0)=0\) and the chosen \(K\)-germ is
  holomorphic at \(Y=0\);
\item
  \(U=2\) is removable by (H3.5) and \[
  \sqrt{1-F(U)/4}=\pm\frac{U-2}{2}
  \] locally;
\item
  no claim about regularity at \(U=4\) is required.
\end{itemize}

Thus the chosen germ \(g\) analytically continues holomorphically
throughout \[
\boxed{
\Omega_U=\mathbf C\setminus[4,\infty).
}
\tag{H4.2}
\] Indeed, the only potential interior singularities of the pulled-back
equation in this simply connected domain are \(0\) and \(2\), and both
are removable for \(g\).

This is the precise continuation statement needed below.

\begin{center}\rule{0.5\linewidth}{0.5pt}\end{center}

\subsubsection{H5. Base and fold maps}\label{h5.-base-and-fold-maps}

Define \[
\boxed{
\phi_b(z)=\frac{16z}{(1+z)^2}.
}
\] This uniformizes the slit plane \(\Omega_U\). One convenient
verification is to put \[
q=\frac{1-z}{1+z}.
\] For \(|z|<1\), \(\Re q>0\), and \[
\phi_b(z)=4(1-q^2).
\] Since squaring maps the right half-plane biholomorphically onto
\(\mathbf C\setminus(-\infty,0]\), \[
\boxed{
\phi_b(\mathbb D)=\mathbf C\setminus[4,\infty).
}
\tag{H5.1}
\]

Now set \[
\phi_f=F\circ\phi_b.
\] A direct calculation gives \[
\boxed{
\phi_f(z)
=
\frac{64z(1-z)^2}{(1+z)^4}.
}
\tag{H5.2}
\] Also \[
\boxed{
\phi_b'(0)=16,
\qquad
\phi_f'(0)=64.
}
\tag{H5.3}
\]

Because \(g\) is holomorphic on \(\Omega_U\), (H5.1) implies \[
g(\phi_b(z))
\] is holomorphic for \(|z|<1\). As germs at \(z=0\), \[
g(\phi_b(z))
=
K(F(\phi_b(z)))
=
K(\phi_f(z)).
\] Hence, with the right-hand side interpreted by analytic continuation
from the germ at \(z=0\), \[
\boxed{
K(\phi_f(z))
=
g(\phi_b(z))
\quad\text{is holomorphic for }|z|<1.
}
\tag{H5.4}
\]

This proves that the fold is carried by the actual conditional
\(K\)-channel.

\begin{center}\rule{0.5\linewidth}{0.5pt}\end{center}

\subsubsection{H6. Coordinate and normalization
warning}\label{h6.-coordinate-and-normalization-warning}

The construction above must not be confused with the older fixed-fold
modular quotient.

Here \[
U\longmapsto Y=F(U)=4U-U^2
\] is introduced only as an auxiliary quadratic cover that removes the
square-root monodromy of the \textbf{already constructed
reflection-channel function \(K(Y)\)}.

No assertion is made that \(U\) is a modular quotient coordinate, nor
that \(F\) identifies the reflection quotient with the older fixed-fold
quotient.

Accordingly, the previously identified common-coordinate obstruction for
the old fixed-fold construction is avoided: the present statement is
simply the analytic pullback \[
K(Y)\longmapsto K(F(U))
\] inside the same \(Y\)-channel.

There is also no \(Y\mapsto Y/4\) rescaling here. For the unscaled
reflection coordinate the base slit map itself is \[
Y=\frac{16z}{(1+z)^2},
\] with derivative \(16\) at the origin. Therefore the value
\(\phi_f'(0)=64\) must be compared with the unscaled base derivative
\(16\), not with the obsolete derivative-\(4\) normalization.

This proves that the displayed fold map is a valid analytic map for the
\(K\)-channel. No assertion of uniqueness or maximality of the
continuation map is needed.

\begin{center}\rule{0.5\linewidth}{0.5pt}\end{center}

\section{Section I --- Mixed Bost energies and evaluation
heights}\label{section-i-mixed-bost-energies-and-evaluation-heights}

\subsubsection{I1. Uniformization and exact first
contact}\label{i1.-uniformization-and-exact-first-contact}

Put \[
q=\frac{1-z}{1+z}.
\] The unit disk is mapped by \(q\) to the right half-plane, and \[
\phi_b(z)=4(1-q^2).
\] Since squaring maps the right half-plane biholomorphically onto
\(\mathbf C\setminus(-\infty,0]\), it follows that \[
\phi_b(\mathbb D)=\mathbf C\setminus[4,\infty).
\] On \(\mathbb T\), \(q\) is purely imaginary, so
\(\phi_b(\mathbb T)=[4,\infty]\).

For the fold, define \[
a(t)=\frac{4t}{(1-t)^2}.
\] A direct calculation gives the stronger identity \[
\boxed{\phi_f(t)=\phi_b(a(t)).}
\] Indeed, \[
\phi_b\!\left(\frac{4t}{(1-t)^2}\right)
=\frac{64t(1-t)^2}{(1+t)^4}.
\] For \(|t|=r<1\), \[
|a(t)|\le \frac{4r}{(1-r)^2},
\] with equality at the positive real point \(t=r\). Hence the fold
image stays in the base slit complement exactly while \[
\frac{4r}{(1-r)^2}<1.
\] First contact therefore occurs at the unique root in \((0,1)\) of \[
4r=(1-r)^2,
\] namely \[
\boxed{r_0=3-2\sqrt2=(\sqrt2-1)^2.}
\] At that point \(a(r_0)=1\), so \[
\phi_f(r_0)=\phi_b(1)=4.
\] This proves that the contact at the positive real point is genuinely
the first contact; there is no earlier off-axis contact.

Since \[
r_0=(1+\sqrt2)^{-2},
\] we obtain on the positive real branch \[
L=-\log r_0
=2\log(1+\sqrt2)
=\log(3+2\sqrt2)
=\boxed{\operatorname{arcosh}3}.
\]

Equivalently, eliminating \(a\) gives \[
16r(1-r)^2=(1+r)^4,
\] whose reciprocal quartic reduces after \(u=r+r^{-1}\) to
\((u-6)^2=0\).

\subsubsection{I2. Exact first-contact mixed
energy}\label{i2.-exact-first-contact-mixed-energy}

Define \[
T(r)=\iint_{\mathbb T^2}
\log|\phi_b(z)-\phi_f(rw)|\,d\mu(z)d\mu(w).
\] For \(|a|\le1\), the exact factorization \[
\phi_b(z)-\phi_b(a)
=
\frac{16(z-a)(1-az)}{(1+z)^2(1+a)^2}
\] and Jensen's/Haar mean identities give \[
\int_{\mathbb T}
\log|\phi_b(z)-\phi_b(a)|\,d\mu(z)
=\log16-2\log|1+a|.
\tag{I.1}
\] (The boundary case \(|a|=1\) is obtained directly by the same
integrable logarithmic means.)

For \(0\le r\le r_0\), put \(t=rw\) and use \(\phi_f(t)=\phi_b(a(t))\).
Moreover \[
1+a(t)
=1+\frac{4t}{(1-t)^2}
=\frac{(1+t)^2}{(1-t)^2}.
\] Therefore \[
\int_{\mathbb T}\log|1+a(rw)|\,d\mu(w)
=2\int_{\mathbb T}\log|1+rw|\,d\mu(w)
-2\int_{\mathbb T}\log|1-rw|\,d\mu(w)
=0,
\] because \(r<1\). Averaging (I.1) gives the stronger statement \[
\boxed{T(r)=\log16\qquad(0\le r\le r_0).}
\] In particular, \[
\boxed{T(r_0)=\log16=4\log2.}
\]

\subsubsection{I3. Exact endpoint
energy}\label{i3.-exact-endpoint-energy}

Write \[
z=e^{2ia},\qquad w=e^{2ib},
\] with \(a,b\) uniform modulo \(\pi\). Then \[
\phi_b(e^{2ia})=4\sec^2a,
\] and \[
\phi_f(e^{2ib})=-16\tan^2b\,\sec^2b.
\] Let \[
X=\tan a,\qquad Y=\tan b.
\] Then \(X,Y\) are independent standard Cauchy variables, so \[
T(1)
=\log4+
\mathbf E\log\!\left(1+X^2+4Y^2+4Y^4\right).
\] Since \[
1+4Y^2+4Y^4=(1+2Y^2)^2,
\] this is \[
T(1)
=\log4+
\mathbf E\log\!\left(X^2+(1+2Y^2)^2\right).
\] For \(A>0\), \[
\frac1\pi\int_{-\infty}^{\infty}
\frac{\log(x^2+A^2)}{1+x^2}\,dx
=2\log(1+A).
\tag{I.2}
\] Applying (I.2) first in \(X\), with \(A=1+2Y^2\), gives \[
\mathbf E_X\log\!\left(X^2+(1+2Y^2)^2\right)
=2\log(2+2Y^2)
=2\log2+2\log(1+Y^2).
\] Applying (I.2) again with \(A=1\) gives \[
\mathbf E\log(1+Y^2)=2\log2.
\] Hence \[
\boxed{T(1)=8\log2.}
\]

For the base self-energy, use the same factorization with
\(a=w\in\mathbb T\). Averaging first in \(z\) and then in \(w\), \[
\iint_{\mathbb T^2}
\log|\phi_b(z)-\phi_b(w)|\,d\mu(z)d\mu(w)
=\log16,
\] so \[
\boxed{\overline L_b^{\,2}=\log16.}
\]

\subsubsection{I4. Mixed Poincaré--Lelong
formula}\label{i4.-mixed-poincaruxe9lelong-formula}

Let
\(\alpha,\beta:(\overline{\mathbb D},0)\to(\mathbf P^1(\mathbf C),0)\)
be holomorphic as projective maps, with \(\alpha'(0),\beta'(0)\ne0\),
and put \[
\beta_r(w)=\beta(rw),\qquad 0<r\le1.
\] Let \(\overline L_\alpha\) and \(\overline L_{\beta,r}\) be the
Bost--Charles metrics on \(O(1)\) induced by \(\alpha\) and \(\beta_r\).
Then \[
\boxed{
\overline L_\alpha\cdot\overline L_{\beta,r}
=
\iint_{\mathbb T^2}
\log|\alpha(z)-\beta(rw)|\,d\mu(z)d\mu(w).
}
\tag{I.3}
\]

Here is the two-map proof. With \[
dd^c=\frac{i}{\pi}\partial\bar\partial,
\qquad
g(w)=\log^+|w|^{-1},
\] one has \[
dd^c g=\mu-\delta_0,
\qquad
c_1(\overline L_\alpha)=\alpha_*\mu.
\] For fixed \(z\in\mathbb T\), Poincaré--Lelong gives \[
\beta_r^*\delta_{\alpha(z)}
=dd^c\log|\beta_r(w)-\alpha(z)|.
\] Green--Stokes therefore yields \[
\int_{\mathbb D}g(w)\,\beta_r^*\delta_{\alpha(z)}
=
\int_{\mathbb T}\log|\beta(rw)-\alpha(z)|\,d\mu(w)
-\log|\alpha(z)|,
\] because \(\beta_r(0)=0\). Integrating in \(z\), and using
Poisson--Jensen for the fiber degree term \[
\widehat{\deg}(\overline L_\alpha|_{0})
=\int_{\mathbb T}\log|\alpha(z)|\,d\mu(z),
\] the last term cancels exactly. This gives (I.3). The proof is
bilinear and nowhere requires \(\alpha=\beta\).

The possible poles of the displayed rational maps on \(\mathbb T\) are
harmless here: as maps to \(\mathbf P^1\) they are holomorphic, and the
resulting logarithmic boundary singularities are integrable.

\subsubsection{I5. Two-radius evaluation-height
inequality}\label{i5.-two-radius-evaluation-height-inequality}

For a degree-\(D\) source section of source norm at most one whose
pullback through \(\beta_r\) vanishes to exact order \(n\) at the
origin, Poisson--Jensen applied to the normalized pullback gives, under
the usual Bost--Charles admissibility hypotheses, \[
\boxed{
h_\infty(\operatorname{ev}_n;r)
\le
-n\log|r\beta'(0)|
+D\bigl(\overline L_\alpha\cdot\overline L_{\beta,r}\bigr)
+o(D).
}
\tag{I.4}
\] For any fixed finite set of radii, the \(o(D)\) term is uniform; in
particular this applies to the \(O(D)\)-range pivots used in the
filtration.

Now take \[
\alpha=\phi_b,\qquad \beta=\phi_f,
\qquad \beta'(0)=64.
\] At \(r=1\), (I.4) gives \[
h_{\infty,f}(n)
\le -n\log64+DT(1)+o(D).
\] At \(r=r_0\), using \(-\log r_0=L\), \[
h_{\infty,f}(n)
\le -n\log64+nL+DT(r_0)+o(D).
\] Taking the better of the two bounds gives \[
\boxed{
h_{\infty,f}(n)
\le
-n\log64
+DT(1)
-\bigl(D\Delta-nL\bigr)_+
+o(D),
}
\] where \[
\Delta=T(1)-T(r_0)=4\log2.
\]

\section{Section J --- Half-fold
optimization}\label{section-j-half-fold-optimization}

\subsubsection{J1. Normalized pivot order and the discrete fold
set}\label{j1.-normalized-pivot-order-and-the-discrete-fold-set}

Put \[
s=\frac nD.
\] Since the first \(2D\) coefficient evaluations are independent, the
pivot orders lie in \[
0\le n\le 2D-1,
\qquad
0\le s<2.
\]

After the endpoint/source-intersection terms are separated in the global
ledger, the order-dependent gain obtained by treating a pivot of order
\(n\) as a fold pivot rather than a base pivot is \[
 n\log\frac{64}{16}
 +(D\Delta-nL)_+
 +o(D).
\] Thus, with \[
d=\log4,
\] this is \[
\boxed{
Dg(s)+o(D),
\qquad
g(s)=ds+(\Delta-Ls)_+.
}
\]

Let \[
S_D\subset\{0,1,\dots,2D-1\}
\] be the pivot orders assigned to fold directions. Since there are
exactly \(D\) fold source directions, \[
|S_D|=D.
\] No assumption on how these \(D\) indices are distributed is made.

\subsubsection{J2. Exact discrete-to-continuous
reduction}\label{j2.-exact-discrete-to-continuous-reduction}

For each \(n\in S_D\), thicken the normalized grid point \(n/D\) to the
cell \[
\left[\frac nD,\frac{n+1}{D}\right),
\] and set \[
F_D=
\bigcup_{n\in S_D}
\left[\frac nD,\frac{n+1}{D}\right)
\subset[0,2].
\] Then \[
|F_D|=1
\] exactly.

The function \(g\) is Lipschitz on \([0,2]\). Hence, uniformly in the
choice of \(S_D\), \[
\int_{F_D}g(s)\,ds
=
\frac1D\sum_{n\in S_D}g\!\left(\frac nD\right)
+O\!\left(\frac1D\right).
\] Equivalently, \[
D\sum_{n\in S_D}g\!\left(\frac nD\right)
=
D^2\int_{F_D}g(s)\,ds+O(D).
\] Therefore it is enough to minimize \[
\int_Fg(s)\,ds
\] over measurable \(F\subset[0,2]\) with \(|F|=1\).

This step removes any need to assume that the discrete fold sets possess
a limiting density or a canonical limiting distribution.

\subsubsection{J3. Rearrangement
minimization}\label{j3.-rearrangement-minimization}

Let \[
c=\frac{\Delta}{L}.
\] For the present constants, \[
d=2\log2,
\qquad
\Delta=4\log2=2d.
\] Moreover \[
d<L,
\] because \[
e^d=4<3+2\sqrt2=e^L.
\] Thus \[
0<\frac dL<1.
\]

The function \(g\) is affine with slope \[
d-L<0
\] on \([0,c]\), and affine with slope \[
d>0
\] on \([c,2]\). Hence \(g\) is strictly V-shaped and its sublevel sets
are intervals.

By the standard bathtub/rearrangement principle, a measurable set of
fixed measure minimizing the integral of \(g\) is a sublevel set of
\(g\). Since \(g\) is strictly monotone on the two sides of its unique
minimum, the measure-one minimizer is an interval \[
F_*=[a,a+1]
\] whose endpoints have equal \(g\)-value.

Since the left endpoint lies before the kink and the right endpoint
after it, \[
\Delta+(d-L)a=d(a+1).
\] Therefore \[
\boxed{
a=\frac{\Delta-d}{L}
=\frac{2\log2}{\operatorname{arcosh}3}.
}
\] Because \(\Delta-d=d\) and \(0<d/L<1\), one has \[
0<a<1,
\] so \(F_*\subset[0,2]\). Also \[
c-a=\frac dL>0,
\qquad
a+1-c=1-\frac dL>0,
\] so the kink indeed lies strictly inside the minimizing interval.

\subsubsection{J4. Exact minimum}\label{j4.-exact-minimum}

Write again \(c=\Delta/L\). Then \[
\int_{a}^{a+1}g(s)\,ds
=
\int_a^{a+1} d\,s\,ds
+
\int_a^c(\Delta-Ls)\,ds.
\] The first term is \[
d\left(a+\frac12\right),
\] and the second is the area of a triangle of base \(d/L\) and height
\(d\): \[
\frac{d^2}{2L}.
\] Consequently \[
I_*
=
\frac d2+
\frac{d\Delta-d^2/2}{L}.
\] Substituting \(d=2\log2\) and \(\Delta=4\log2\) gives \[
\boxed{
I_*
=
\log2+
\frac{6(\log2)^2}{\operatorname{arcosh}3}.
}
\] Numerically, \[
I_*\approx2.328502565588796.
\]

\subsubsection{J5. Discrete total bound}\label{j5.-discrete-total-bound}

For every assignment of the \(D\) fold source directions to \(D\) of the
first \(2D\) Hermite pivot orders, \[
\begin{aligned}
\text{total fold-flag gain}
&\ge
D\sum_{n\in S_D}
 g\!\left(\frac nD\right)
+o(D^2)\\
&\ge
I_*D^2+o(D^2).
\end{aligned}
\] Here the first \(o(D^2)\) follows from the uniform \(o(D)\) per-pivot
remainder imported from Section I (or from any equivalent aggregate
\(o(D^2)\) estimate).

Hence \[
\boxed{
\text{total fold flag contributes at least }
I_*D^2+o(D^2),
\qquad
I_*=\log2+
\frac{6(\log2)^2}{\operatorname{arcosh}3}.
}
\]

No conjecture about a canonical or limiting pivot distribution is
required.

\section{Section K-Global --- Adelic
assembly}\label{section-k-global-adelic-assembly}

\subsection{K8. Odd-prime denominator
cost}\label{k8.-odd-prime-denominator-cost}

Section G works only at odd finite primes and gives \[
\boxed{
C_{\rm odd}=3.
}
\tag{K8.1}
\]

Prime \(2\) is excluded, so there is no overlap with K7.

\begin{center}\rule{0.5\linewidth}{0.5pt}\end{center}

\subsection{K9. Archimedean constants}\label{k9.-archimedean-constants}

Section I gives \[
\boxed{
b_\infty=\log16,
\qquad
T(1)=8\log2.
}
\tag{K9.1}
\]

Section J gives \[
\boxed{
I_*
=
\log2+
\frac{6(\log2)^2}{\operatorname{arcosh}3}.
}
\tag{K9.2}
\]

All pivots satisfy \[
0\le n<2D,
\] so every CD normalization estimate requiring \(n=O(D)\) is
\(o(D^2)\).

The global basis change \[
(P,R_0)\longleftrightarrow(P,K)
\] has determinant \(1\), so the Archimedean base/fold decomposition is
on the same rational determinant line.

\begin{center}\rule{0.5\linewidth}{0.5pt}\end{center}

\subsection{K10. Complete determinant/basis-change
ledger}\label{k10.-complete-determinantbasis-change-ledger}

\begin{longtable}[]{@{}
  >{\raggedright\arraybackslash}p{(\columnwidth - 6\tabcolsep) * \real{0.2000}}
  >{\raggedleft\arraybackslash}p{(\columnwidth - 6\tabcolsep) * \real{0.2667}}
  >{\centering\arraybackslash}p{(\columnwidth - 6\tabcolsep) * \real{0.3333}}
  >{\raggedright\arraybackslash}p{(\columnwidth - 6\tabcolsep) * \real{0.2000}}@{}}
\toprule\noalign{}
\begin{minipage}[b]{\linewidth}\raggedright
Transition
\end{minipage} & \begin{minipage}[b]{\linewidth}\raggedleft
Determinant / size
\end{minipage} & \begin{minipage}[b]{\linewidth}\centering
Quadratic?
\end{minipage} & \begin{minipage}[b]{\linewidth}\raggedright
Charged where
\end{minipage} \\
\midrule\noalign{}
\endhead
\bottomrule\noalign{}
\endlastfoot
Coinvariants \(\to\) invariants & \(2^D\) & No & D \\
Integral rationality-adapted basis & \((2b)^D\) & No & D \\
\((P,R_0)\to(P,K)\) & \(1\) & No & global \\
Divide jets by \(P(64z)\) & determinant \(1\) & No & E/K3 \\
\(\zeta\mapsto\zeta(1-4\zeta)\) on full jets & determinant \(1\) & No &
K2 \\
Column reversal & sign only & No & E \\
CD unit matrices & units & No & K4 \\
Literal CD \(W_n\)-normalization & \(O(\log D)\) & No & B/K4 \\
\(Y=64z\) Hermite scaling & explicit in K5.3 & \textbf{Yes} & K5 only \\
Source transition \(S\) & \(v_2\det S=D^2+2S_D\) & \textbf{Yes} & K6 \\
Target transition \(T\) & \(-v_2\det T=D^2-5D+2S_D\) & \textbf{Yes} &
K6 \\
Canonical reflection Smith mass & \(R_D\) & \textbf{Yes} & C \\
Third reflection-scale reserve & \(R_D\) & \textbf{Reserve} & K7 \\
Odd-prime denominator lattice & \(3D^2+o(D^2)\) & \textbf{Yes} & G,
\(p\ne2\) \\
\end{longtable}

The first two \(R_D\)-scale contributions are the actual source and
target compatibility transitions. The third is deliberately unused
slack.

\begin{center}\rule{0.5\linewidth}{0.5pt}\end{center}

\subsection{K11. Quadratic adelic
ledger}\label{k11.-quadratic-adelic-ledger}

After cancellation of the common all-base source/derivative
contribution, rationality implies \[
\boxed{
C_{2,\mathrm{eff}}
+
2b_\infty
-
C_{\rm odd}
+
I_*
\le
T(1).
}
\tag{K11.1}
\]

Substituting, \[
\boxed{
\log2+2\log16-3+
\log2+
\frac{6(\log2)^2}{\operatorname{arcosh}3}
\le8\log2.
}
\tag{K11.2}
\]

Every omitted contribution is \(o(D^2)\).

Numerically, \[
\mathrm{LHS}\approx5.566827190628304,
\qquad
\mathrm{RHS}\approx5.545177444479562,
\] so the inequality fails by approximately \[
0.021649746148742.
\]

The rigorous final contradiction is delegated to Section L.

\begin{center}\rule{0.5\linewidth}{0.5pt}\end{center}

\section{Section L --- Certified scalar
contradiction}\label{section-l-certified-scalar-contradiction}

\subsubsection{L1. Scalar inequality}\label{l1.-scalar-inequality}

Section K-Global gives, under the assumption \(G\in\mathbf Q\), \[
\log2+2\log16-3+I_*
\le8\log2,
\] where \[
I_*=
\log2+
\frac{6(\log2)^2}{\operatorname{arcosh}3}.
\]

Since \[
2\log16=8\log2,
\] cancelling the common \(8\log2\) terms gives \[
\boxed{
2\log2-3+
\frac{6(\log2)^2}{\operatorname{arcosh}3}
\le0.
}
\]

\subsubsection{L2. Certified positivity}\label{l2.-certified-positivity}

Set \[
\ell=\log2,
\qquad
A=\operatorname{arcosh}3.
\] Define \[
F(\ell,A)=2\ell-3+\frac{6\ell^2}{A}.
\] For \(\ell,A>0\), \[
\frac{\partial F}{\partial\ell}
=
2+\frac{12\ell}{A}>0,
\qquad
\frac{\partial F}{\partial A}
=
-\frac{6\ell^2}{A^2}<0.
\] Thus a lower bound for \(\ell\) and an upper bound for \(A\) give a
lower bound for \(F\).

First, \[
\log2
=
2\operatorname{artanh}\frac13
=
2\sum_{k=0}^{\infty}
\frac1{(2k+1)3^{2k+1}}.
\] Since all terms are positive, \[
\log2
>
2\sum_{k=0}^{5}
\frac1{(2k+1)3^{2k+1}}
=
\frac{15757912}{22733865}.
\] Direct rational comparison gives \[
\frac{15757912}{22733865}
-
\frac{693147}{10^6}
=
\frac{335369}{4546773000000}>0.
\] Hence \[
\boxed{\log2>\frac{693147}{10^6}.}
\]

Next, if \(A=\operatorname{arcosh}3\), then \[
\tanh\frac A2
=
\sqrt{\frac{\cosh A-1}{\cosh A+1}}
=
\frac1{\sqrt2}.
\] Therefore \[
A
=
2\operatorname{artanh}\frac1{\sqrt2}
=
\sqrt2\sum_{k=0}^{\infty}
\frac1{2^k(2k+1)}.
\]

We have \[
\sqrt2<\frac{707107}{500000},
\] since \[
\left(\frac{707107}{500000}\right)^2-2
=
\frac{309449}{250000000000}>0.
\]

Moreover, \[
\sum_{k=12}^{\infty}
\frac1{2^k(2k+1)}
<
\frac1{25}\sum_{k=12}^{\infty}2^{-k}
=
\frac1{51200}.
\] Consequently \[
\begin{aligned}
A
&<
\frac{707107}{500000}
\left(
\sum_{k=0}^{11}\frac1{2^k(2k+1)}
+\frac1{51200}
\right)\\
&=
\frac{215729472659027303}
{122382374400000000}\\
&<
\frac{7051}{4000}.
\end{aligned}
\] Thus \[
\boxed{
\operatorname{arcosh}3<\frac{7051}{4000}.
}
\]

By the monotonicity established above, \[
\begin{aligned}
2\log2-3+
\frac{6(\log2)^2}{\operatorname{arcosh}3}
&>
2\frac{693147}{10^6}-3
+
\frac{
6(693147/10^6)^2
}{
7051/4000
}\\
&=
\frac{19078165077}{881375000000}.
\end{aligned}
\] Finally, \[
\frac{19078165077}{881375000000}
-
\frac{27}{1250}
=
\frac{40465077}{881375000000}>0.
\] Hence \[
\boxed{
2\log2-3+
\frac{6(\log2)^2}{\operatorname{arcosh}3}
>
\frac{27}{1250}
=0.0216>0.
}
\]

This contradicts L1.

Therefore the assumption \(G\in\mathbf Q\) is impossible.

\subsubsection{Exported conclusion}\label{exported-conclusion}

\[
\boxed{G=\beta(2)\notin\mathbf Q.}
\]

\section{Proof of the main theorem}\label{proof-of-the-main-theorem}

Assume for contradiction that \(G\in\mathbf Q\). The exact arithmetic
theorem of Section A gives \[
\Omega_D=4D^2+4D-4S_D,
\qquad
S_D=\sum_{m=1}^D s_2(m).
\] The common-source comparison of Section K-Finite shows that the two
actual source/target compatibility transitions each have one-sided size
\(R_D+O(D\log D)\), where Section C gives \[
R_D=D^2-\lceil D/2\rceil.
\] Retaining one further complete \(R_D\) as a conservative reserve
leaves \[
(\Omega_D-3R_D)\log2
=\log2\,D^2+o(D^2).
\] Thus the effective \(2\)-adic coefficient is \(\log2\).

Section G proves directly in the final \(Y\)-coordinate that the
odd-prime source has denominator profile \((0,2)\), whose rearranged
cost is \(3D^2+o(D^2)\). Sections H and I construct the actual
conditional fold and establish \[
\phi_b'(0)=16,
\qquad
\phi_f'(0)=64,
\qquad
T(r_0)=4\log2,
\qquad
T(1)=8\log2,
\] with \[
r_0=3-2\sqrt2,
\qquad
L=-\log r_0=\operatorname{arcosh}3.
\] Section J proves, without assuming any limiting pivot distribution,
that the \(D\) fold directions among the \(2D\) independent Hermite
pivots contribute at least \[
I_*D^2+o(D^2),
\qquad
I_*=\log2+\frac{6(\log2)^2}{\operatorname{arcosh}3}.
\]

The determinant-line/product-formula identity of Section K-Global
therefore forces \[
\log2+2\log16-3+I_*\le8\log2.
\] Since \(2\log16=8\log2\), this is \[
2\log2-3+\frac{6(\log2)^2}{\operatorname{arcosh}3}\le0.
\] Section L proves by rational interval estimates that \[
2\log2-3+\frac{6(\log2)^2}{\operatorname{arcosh}3}
>\frac{27}{1250}>0,
\] a contradiction. Hence \(G\notin\mathbf Q\). \(\square\)

\section{Remarks on verification
status}\label{remarks-on-verification-status}

The assembled text above is intentionally more conservative than the
sharpest intermediate estimates. In particular, it retains three
complete reflection-scale masses although the explicit compatibility
theorem exhibits only two actual quadratic source/target transitions. It
also keeps the prime \(2\) separate from the odd-prime denominator
asymptotic, charges the full \(Y=64z\) coordinate factor exactly once,
and uses a distribution-free fold optimization.

The source modules supplied for this manuscript were returned as
individually audited/corrected units. Because Theorem \ref{thm:main}
would resolve a classical open problem, the present manuscript should
nevertheless be independently checked as a whole before being treated as
an established result. In particular, the most important global
verification target is the exact determinant-line factorization in
Section K, since that is where the independently proved local modules
are joined.

\section{References}\label{references}

\begin{enumerate}
\def\labelenumi{\arabic{enumi}.}
\tightlist
\item
  F. Beukers, \emph{Irrationality of some p-adic L-values},
  arXiv:math/0603277.
\item
  T. Rivoal, \emph{Nombres d'Euler, approximants de Padé et constante de
  Catalan}, Ramanujan Journal \textbf{11} (2006), 199--214.
\item
  T. Rivoal and W. Zudilin, \emph{Diophantine properties of numbers
  related to Catalan's constant}, Mathematische Annalen \textbf{326}
  (2003), 705--721.
\item
  W. Zudilin, \emph{An Apéry-like difference equation for Catalan's
  constant}, Electronic Journal of Combinatorics \textbf{10} (2003),
  R14.
\item
  Yu. V. Nesterenko, \emph{On Catalan's Constant}, Proceedings of the
  Steklov Institute of Mathematics \textbf{292} (2016), 153--170.
\item
  F. Calegari, V. Dimitrov and Y. Tang, \emph{The linear independence of
  1, \(\zeta(2)\), and \(L(2,\chi_{-3})\)}, arXiv:2408.15403.
\item
  J.-B. Bost and F. Charles, \emph{Quasi-projective and formal-analytic
  arithmetic surfaces}, arXiv:2206.14242.
\end{enumerate}

\end{document}
